% !TEX root = trayguard_paper.tex
\documentclass[10pt,a4paper]{article}

% ── Packages ───────────────────────────────────────────────────────────
\usepackage[utf8]{inputenc}
\usepackage[T1]{fontenc}
\usepackage{geometry}
\geometry{margin=1in}
\usepackage{setspace}
\setstretch{1.15}
\usepackage{parskip}
\usepackage{hyperref}
\hypersetup{
    colorlinks=true,
    linkcolor=blue,
    citecolor=blue,
    urlcolor=blue,
}
\usepackage{booktabs}
\usepackage{longtable}
\usepackage{array}
\usepackage{tabularx}
\usepackage{verbatim}
\usepackage{amsmath}
\usepackage{amssymb}
\usepackage{graphicx}
\usepackage{float}
\usepackage{xcolor}
\usepackage{listings}
\lstset{basicstyle=\footnotesize\ttfamily,breaklines=true,columns=fullflexible}

\setcounter{tocdepth}{3}
\setcounter{secnumdepth}{3}

\title{TrayGuard\\A Local Tray Training And Assessment System}
\author{}
\date{}

\begin{document}

\maketitle
\tableofcontents
\newpage

\section{Abstract / Executive Summary}
\label{sec:1-abstract-executive-summary}
TODO: Write this section last after the body of the report is complete. It
should summarize the final problem, the selected TrayGuard direction, the
strongest evidence, the main design decisions, the current limitations, and the
recommended future validation path.
\section{Introduction And Project Overview}
\label{sec:2--introduction-and-project-overview}
Tray content errors during surgical instrument reconstruction are recurring
operational failures that can delay surgery or compromise patient safety. The
workforce faces structural pressure from growing demand, high turnover, and
limited training capacity. Current training blends online coursework,
classroom instruction, and hands-on externships that are capacity-constrained by
clinical-site availability, preceptor attention, and local
teaching materials.

TrayGuard originally targeted a computer-vision-based tray checker: a camera
above the assembly workstation would detect each instrument as the technician
placed it, compare detections against the tray list, and return a go or no-go
verdict before sealing. The CV approach failed early because it required the
project to validate detection accuracy, lighting tolerance, occluded-item
handling, and workflow integration inside a live SPD environment, resources
not available within our project scope. The project pivoted to a training-centered
alternative that preserves the same error categories (missing, wrong, extra,
misidentified, and miscounted items) but measures them in a simulated local tray
task using AprilTag instrument cards, retrieval study, and feedback modes.

This Functional Design Review serves as the literature-grounded design handoff
for a future implementation team. It defines the problem, justifies the
training-centered solution, specifies requirements and system architecture,
provides a study plan for novice evaluation, and documents the scope boundaries
that prevent overclaiming.
\section{Problem Definition And Scope}
\label{sec:3--problem-definition-and-scope}
\subsection{Precise Problem Statement}
\label{sec:3-1-precise-problem-statement}
Sterile processing technicians receive used instruments from the OR, clean,
inspect, assemble procedure-specific trays, package, sterilize, store, and
distribute sets for surgical use. Mistakes can occur at several points in that
workflow: a technician may confuse two visually similar instruments, place the
wrong size or pattern into a tray, omit or add an instrument, use the wrong
count sheet, or fail to catch a documentation problem. In the OR, those
mistakes appear as missing, wrong, extra, damaged, contaminated, or otherwise
unavailable instruments.

Published rates vary by hospital, measurement method, and error definition, but
the available evidence suggests that these are recurring operational failures,
not isolated accidents. Alfred et al. analyzed 3,900 tray defects across 41,799
cases, or roughly 9 defects per 100 cases, and found that 55.0\% of recorded
defects occurred during assembly
\cite{alfred-et-al-2021}. Zhu et al. found
398 packaging errors among 33,839 surgical instrument packages, or about 1.2\%
of packages, including wrong specifications, incomplete packages, and missing
instruments \cite{zhu-et-al-2019}. Nichol et al.
observed 236 surgical instrument errors affecting 147 cases, with most errors
tied to visualization work such as inspection, identification, function
checking, and tray assembly
\cite{nichol-et-al-2024}.

The general problem is therefore:
\begin{quote}
Sterile processing workflows do not always reliably produce complete,
correct, functional, sterile, and ready-for-use surgical instrument sets.
\end{quote}
\subsection{Practice Context}
\label{sec:3-2-practice-context}
The problem appears across the reusable stainless surgical instrument chain. A simplified
SPD-to-OR cycle includes the following stages
\cite{george-et-al-2024}:
\begin{enumerate}
  \item Point-of-use preparation: at the end of a procedure, instruments are prepared for transfer from the OR to sterile processing. This can include gross cleaning, use of sterile water, and enzymatic pretreatment to keep material from drying on instruments.
  \item Transport to sterile processing: used instruments are moved from the OR or procedural area to the sterile processing unit.
  \item Decontamination: instruments are cleaned and decontaminated. George et al. describe manual scrubbing, ultrasonic cleaning, and automated washer systems as common parts of this stage.
  \item Assembly and inspection: after decontamination, instruments move to the assembly area. Technicians inspect instruments, assemble instrument pans or sets, replace missing items, and place instruments into containers or wraps for sterilization.
  \item Sterilization: packaged instruments are sterilized using methods appropriate to the device and material. George et al. describe high-temperature methods such as autoclave sterilization and low-temperature methods such as chemical or radiation-based sterilization.
  \item Storage: once sterilized, instruments and trays are cooled as needed and placed in storage until they are required again.
  \item Return to use: when needed for surgery, stored instruments are obtained and brought back to the OR for the next case.
\end{enumerate}
Nichol and Saari's risk model provides a more detailed version of this cycle:
for simple instruments, they mapped 104 human tasks and found that 91\% of
modeled error risk occurred during the sterile-processing portion of the cycle,
especially the decontamination and assembly-and-inspection steps
\cite{nichol-and-saari-risk-modeling-2023}.
Those steps are high risk because they contain the most complicated,
non-routine tasks: instruments must be separated, cleaned, inspected for
bioburden and damage, identified, sorted, and rebuilt into the correct tray, and
all eight highest-risk tasks in the model involved human visualization,
inspection, identification, or sorting.
\subsection{Narrowed Problem}
\label{sec:3-3-narrowed-problem}
Within the full SPD-to-OR cycle, this project narrows to tray reconstruction
and verification: the point where cleaned instruments are inspected,
identified, matched to count sheets, and rebuilt into procedure-specific trays.
This is the most defensible focus because the highest-risk parts of the cycle
depend on human visualization, sorting, and local tray knowledge.
This focus also addresses a likely stakeholder objection. Even if a tray comes
from a specific surgery and is expected to return with the same contents, it
does not remain a sealed, unchanged kit during reprocessing. Instruments are
transported to decontamination, opened, separated, disassembled when required,
manually cleaned or transferred onto washer racks, and exposed to mechanical
washing before moving to the clean preparation area
\cite{steris-cleaning-instruments-accessed-2026},
\cite{cdc},
\cite{public-health-ontario-2013},
\cite{medline-transport-2026}. After cleaning,
technicians manually rebuild sets against tray-specific count sheets, creating
opportunities for items to be mixed between sets, left behind, substituted from
backup inventory, or misidentified because of incomplete count-sheet data,
local jargon, or look-alike instruments
\cite{purdue-count-sheets-accessed-2026},
\cite{alfred-et-al-2021},
\cite{nadeau-2024}.
\subsection{Stakeholders And Importance}
\label{sec:3-4-stakeholders-and-importance}
Instrument readiness matters because surgery is time-sensitive,
coordination-heavy work. When an instrument set is wrong, incomplete, unclean,
damaged, or unavailable, the OR team may need to pause, search for replacement
instruments, open backup trays, change the procedure flow, or delay the case
\cite{nichol-et-al-2024},
\cite{pennsylvania-patient-safety-authority-2006}.
Dirty or damaged instruments add a direct safety concern. A 2024 hospital
inspection reported pitting, stains, sticky residue, rust, scratches, and other
concerns in randomly selected ready-for-use trays
\cite{association-of-health-care-journalists-2024}.
The Pennsylvania Patient Safety Authority warns that contaminated instruments
can put patients at risk of surgical site infection and can also cause lost OR
time if discovered after a procedure has begun
\cite{pennsylvania-patient-safety-authority-2006}.
The stakeholders include:
\begin{itemize}
  \item Patients, who depend on sterile, functional instruments during surgery.
  \item Surgeons and OR teams, who need complete and usable sets to maintain procedure flow.
  \item SPD technicians, who perform detailed, high-volume work under production pressure.
  \item SPD supervisors and educators, who manage training, quality assurance, staffing, documentation, and feedback loops.
  \item Hospitals, which lose time and capacity when instrument issues delay cases or require rework.
\end{itemize}
\subsection{General Interest Beyond One Local Practice}
\label{sec:3-5-general-interest-beyond-one-local-practice}
This problem is broadly relevant because reusable surgical instruments, SPDs,
procedure-specific trays, count sheets, sterilization records, and OR
dependence on prepared sets are common across hospitals. The exact instruments,
tray lists, vendor systems, and surgeon preferences vary by site, while the
readiness challenge repeats across settings: hospitals must transform complex,
locally defined instrument requirements into complete and ready surgical sets.
The local tray-reconstruction problem connects to broader work on surgical
tray optimization, inventory visibility, and instrument-set readiness. dos
Santos et al. describe tray rationalization as deciding which instruments and
quantities belong in which trays and which trays are needed for procedures
\cite{dos-santos-et-al-2021}. A 2025
industry benchmark report summarized by Healthcare Purchasing News found that
58\% of surveyed SPD and OR leaders reported surgery delays caused by instrument
sets that were missing, incomplete, or not ready on time
\cite{hpn-2025}.
The same benchmark report emphasized human error, inventory gaps, limited data
visibility, and real-time visibility into tray and instrument location as major
readiness concerns
\cite{aesculap-and-ascendco-health-2025}.
This makes the problem generalizable without making it generic: each hospital
has its own instruments and tray rules, but many hospitals face the same
underlying need to verify that a locally defined set is complete, correct,
clean, functional, and available before surgery.
\section{Root Cause Analysis}
\label{sec:4--root-cause-analysis}
We have already narrowed the problem to a specific point in the sterile
processing workflow. The fishbone head for the RCA is therefore:
\begin{quote}
Tray content errors during tray reconstruction.
\end{quote}
This effect includes missing instruments, wrong instruments, extra instruments,
wrong quantities, visible damage or function problems, and item selections that
conflict with the count sheet or local tray rule
\cite{alfred-et-al-2021};
\cite{nichol-et-al-2024};
\cite{zhu-et-al-2019}.
\subsection{Fishbone  Summary}
\label{sec:4-1-fishbone-summary}
\begin{itemize}
  \item \textbf{Product:} Same-family instruments differ only in small features with no labels or markings. Steam autoclaving destroys adhesives and fades inks, laser etching is expensive and unstandardized across manufacturers, and no regulatory body requires SPD-level instrument markings. Instruments are designed for surgical function; SPD identification is a secondary consideration. Wear further hides the differences. Nothing on the tool itself helps tell them apart \cite{alfred-et-al-2021}; \cite{nichol-et-al-2024}.
  \item \textbf{Knowledge:} Three sub-causes converge: (a) knowing the general name is not enough, since workers must see the difference between nearly identical instruments \cite{castillo-gutierrez-2025}; \cite{nichol-et-al-2024}; (b) knowing what a Crile is does not tell you what this tray calls it or how many it expects, since general knowledge and local tray rules are separate domains — the same generic name maps to physically different tools across manufacturers, tray composition is driven by surgeon preference and hospital case mix rather than a universal standard, and local aliases accumulate organically on count sheets with no central naming authority \cite{nadeau-2024}; \cite{dos-santos-et-al-2021}; (c) certification proves a starting point, not local readiness, since this hospital's trays require repeated practice with feedback \cite{hspa-crcst-accessed-2026}; \cite{ofstead-et-al-2023}.
  \item \textbf{Process:} Building a tray forces picking, counting, checking, and arranging into one step with no double-check. Sequential verification would double labor per tray, no accreditation mandate requires a two-person check for tray assembly, and productivity metrics (trays/hour) structurally disincentivize redundant steps. The work uses only eyes and hands. Without aids, small errors pass through \cite{alfred-et-al-2021}; \cite{nichol-et-al-2024}.
  \item \textbf{Information:} When count sheets are incomplete or outdated, workers stop following rules and start guessing \cite{nadeau-2024}. Tray contents change continuously — new instruments, surgeon preference changes, substitutions, vendor loaners — but no dedicated role or budget exists for count-sheet maintenance; the update workflow is manual and low-priority. Missing photos, names, or loaner-tray info adds to the risk \cite{steris-loaner-trays-2021}.
  \item \textbf{Environment:} Pressure and interruptions break focus for detail-heavy visual work. SPD is a cost center whose staffing is structurally constrained while OR capacity — a revenue center — expands independently. Production pressure is built into the workflow: trays must be ready for the next surgery, not when the technician is ready. Short staffing means less training and looser checking \cite{alfred-et-al-2021}; \cite{huang-et-al-2025}. When SPD staffing doesn't keep up with new ORs, backlogs pile up \cite{ingold-2025}.
  \item \textbf{Policy \& Feedback:} Nobody owns the standards — SPD governance is decentralized across shifts and departments; certification bodies (HSPA, CBSPD) recommend but do not mandate recurring competency checks tied to local tray knowledge. Error reporting is manual, slow, and deprioritized; the feedback loop from OR to SPD is broken — errors discovered during surgery are rarely traced back to the specific technician, tray, or training gap. Workers only get tested once at hiring. Slow error reporting lets the same mistakes happen again \cite{loria-2024}; \cite{kovach-2012}; \cite{nichol-et-al-2024}. Training often stops after orientation, and informal mentoring perpetuates inconsistent habits \cite{klacik-2023}.
\end{itemize}
\subsection{Product}\label{sec:4-2-product}
\begin{table}[H]
\centering
\caption{Fishbone: Product root causes}
\label{tab:fishbone-product-root-causes}
\small\renewcommand{\arraystretch}{1.2}\setlength{\tabcolsep}{4pt}
\begin{tabularx}{\textwidth}{|>{\raggedright\arraybackslash}X|>{\raggedright\arraybackslash}X|}
\hline
Observed failure mode & Root cause hypothesis \\
\hline
Instruments within same families (Kelly/Crile, straight/curved Mayo, etc.) differ only in subtle dimensional features — blade angle, tip shape, serration pattern — with no intrinsic labels or physical keying. Wear, staining, and reprocessing cycles further reduce visible distinguishing cues \cite{alfred-et-al-2021}; \cite{lind-2026a}; \cite{nichol-et-al-2024}. & The instruments have no labels, color codes, or markings to tell them apart. Steam autoclaving destroys adhesives and fades most marking materials, laser etching is expensive and unstandardized across manufacturers, and no regulatory body mandates SPD-level instrument identification. Instruments are designed for surgical function with SPD identification as a secondary consideration; the economic incentive for labeling is distributed — benefit accrues to SPD, cost falls on manufacturers \cite{alfred-et-al-2021}; \cite{hspa-instrument-resources-accessed-2026}; \cite{nichol-and-saari-2023}. Workers must memorize subtle details that nothing on the tool itself points out. \\
\hline
\end{tabularx}
\end{table}
\subsection{Knowledge}\label{sec:4-3-knowledge}
\begin{table}[H]
\centering
\caption{Fishbone: Knowledge root causes}
\label{tab:fishbone-knowledge-root-causes}
\small\renewcommand{\arraystretch}{1.2}\setlength{\tabcolsep}{4pt}
\begin{tabularx}{\textwidth}{|>{\raggedright\arraybackslash}X|>{\raggedright\arraybackslash}X|}
\hline
Observed failure mode & Root cause hypothesis \\
\hline
Wrong instruments appear in trays and are often similar in type to the intended instrument; instrument identification, inspection, names, uses, and testing points are explicit sterile-processing learning needs \cite{nichol-et-al-2024}; \cite{hspa-instrument-resources-accessed-2026}. & Knowing the general name isn't enough — the technician must see the difference. Kelly and Crile look nearly identical to untrained eyes, and tray errors are predominantly visual \cite{nichol-et-al-2024}; \cite{alfred-et-al-2021}. Sensory skills (sight, touch, hearing, smell) must be developed through hands-on practice and mentorship; they are not automatic \cite{castillo-gutierrez-2025}. The visual gap is structural: a technician who knows that a Kelly is a hemostat with horizontal serrations still cannot distinguish a Kelly from a Crile without practiced perceptual discrimination — verbal knowledge of the instrument class does not transfer to visual identification in context \cite{fast-et-al-2019}; \cite{shreckengost-et-al-2022}. \\
\hline
Tray assembly requires local names, aliases, quantities, placement rules, and count-sheet use; procedure-, surgeon-, and tray-specific decisions determine which instruments and quantities belong in a set \cite{nadeau-2024}; \cite{dos-santos-et-al-2021}; \cite{ahmadi-et-al-2023}. & Knowing what a Crile is doesn't tell you what this tray calls it or how many it expects. General instrument knowledge and local tray rules are separate domains for four reasons that together make a universal instrument list infeasible. First, manufacturers produce physically different tools under the same generic name — a Kelly hemostat from Company A may differ in length, jaw angle, or serration pattern from one from Company B, but both list as "Kelly" on the count sheet \cite{nadeau-2024}. Second, tray composition is an operational document driven by surgeon preference, hospital case mix, and historical accretion — no external authority mandates what goes in each tray \cite{dos-santos-et-al-2021}. Third, local aliases (e.g., "the blue-handle clamp," "Dr. Smith's right-angle") accumulate organically on count sheets with no central nomenclature authority to harmonize them across hospitals \cite{alfred-et-al-2021}. Fourth, count sheets encode tray-specific rules — placement order, quantity per instrument, approved substitutions — that exist only in that hospital's workflow \cite{nadeau-2024}. \\
\hline
Alfred et al. identify technician knowledge, training variation, missing or incorrect photos, and varied names as assembly work-system factors; if names, aliases, and photos are not aligned with the actual tray, users may match the wrong concept to the physical instrument or fail to recognize a local variant \cite{alfred-et-al-2021}. & Tray knowledge is taught on the job without a standard source — no central training curriculum covers this hospital's specific tray configuration, so each preceptor teaches from personal experience and tacit knowledge. Differences across preceptors, shifts, and training tenures produce inconsistent mental models of the same tray among workers on the same line \cite{alfred-et-al-2021}. Structured SPD training programs that use simulation and assessment show that standardizing teaching improves knowledge and operational skills \cite{hu-et-al-2024}; \cite{thurmond-2020}. \\
\hline
HSPA requires 400 hands-on hours for CRCST certification, including 120 hours in preparation and packaging; AORN describes hands-on and simulation training for assembling and wrapping trays; certification commentary frames certification as a baseline requiring ongoing competency checks \cite{hspa-crcst-accessed-2026}; \cite{aorn-staffing-shortage-2025}; \cite{kovach-2012}; \cite{nadeau-2017}. & Certification proves a baseline, not local competence. The CRCST exam tests general instrument knowledge (names, categories, uses) but cannot assess familiarity with this hospital's specific tray configurations, aliases, quantity rules, or surgeon preferences. Those are learned on site after certification \cite{hspa-crcst-accessed-2026}; \cite{ofstead-et-al-2023}. Neither HSPA nor CBSPD requires recertification tied to local tray knowledge, so even certified workers may have unreviewed gaps on this hospital's current trays \cite{kovach-2012}; \cite{nadeau-2017}. \\
\hline
\end{tabularx}
\end{table}
\subsection{Process}\label{sec:4-4-process}
\begin{table}[H]
\centering
\caption{Fishbone: Process root causes}
\label{tab:fishbone-process-root-causes}
\small\renewcommand{\arraystretch}{1.2}\setlength{\tabcolsep}{4pt}
\begin{tabularx}{\textwidth}{|>{\raggedright\arraybackslash}X|>{\raggedright\arraybackslash}X|}
\hline
Observed failure mode & Root cause hypothesis \\
\hline
Alfred et al. found that 55.0\% of recorded tray defects occurred during assembly, including missing, wrong, damaged, extra, and incorrectly assembled instruments \cite{alfred-et-al-2021}. Instrument damage such as corrosion, stains, burrs, and worn jaw serrations can be missed when inspection relies on rapid visual checks under poor lighting \cite{lind-2026a}. & Building a tray means picking, counting, checking, and arranging items all at once with no double-check system. Sequential verification — one worker picks and arranges, a second worker independently checks — would double labor per tray, and no accreditation mandate (JC, CMS, HSPA) requires a two-person check for tray assembly. Productivity metrics (trays per technician per hour) structurally incentivize single-pass assembly; any verification step that adds time without adding throughput faces organizational resistance \cite{alfred-et-al-2021}; \cite{chen-et-al-2023}. \\
\hline
Nichol et al. found that 88.6\% of observed errors involved visualization tasks such as inspection, identification, function checking, and sorting \cite{nichol-et-al-2024}. & Workers use only their eyes and hands — no magnification, no side-by-side comparison guide, no decision-support tool at the workstation. The assembly station is designed for throughput, not for perceptual support. Without aids that highlight distinguishing features or flag common lookalike pairs, even experienced technicians miss differences under time pressure \cite{nichol-et-al-2024}. \\
\hline
HPN frames following count sheets and inspecting instruments as core preparation, packing, and assembly practices \cite{nadeau-2024}. & Nobody can remember every item in every tray — a single hospital may manage hundreds of tray configurations, each with 30--100+ instruments. The process assumes the worker will cross-reference the count sheet with each selection, but production pressure and cognitive load (simultaneous picking, counting, checking, arranging) cause workers to break from the count sheet and rely on memory, especially for familiar trays \cite{nadeau-2024}. \\
\hline
\end{tabularx}
\end{table}
\subsection{Information}\label{sec:4-5-information}
\begin{table}[H]
\centering
\caption{Fishbone: Information root causes}
\label{tab:fishbone-information-root-causes}
\small\renewcommand{\arraystretch}{1.2}\setlength{\tabcolsep}{4pt}
\begin{tabularx}{\textwidth}{|>{\raggedright\arraybackslash}X|>{\raggedright\arraybackslash}X|}
\hline
Observed failure mode & Root cause hypothesis \\
\hline
Count sheets vary in completeness across shifts and trays; many lack photos, aliases, placement instructions, or current quantities \cite{nadeau-2024}. Missing or outdated identification information, IFUs, and tray listings at the workstation force workers to rely on memory \cite{lind-2025}. & Count sheets become incomplete or outdated because tray contents change continuously — instruments are added, removed, substituted, or replaced as surgeon preferences shift, manufacturers discontinue items, or vendors supply new loaners — but no dedicated role or budget exists for count-sheet maintenance. Updating a count sheet is a manual, discretionary task that competes with production work. Without a formal change-management workflow, sheets accumulate inaccuracies across multiple revision cycles \cite{nadeau-2024}. \\
\hline
Loaner-tray workflows depend on visibility across scheduling, vendor delivery, OR awareness, SPD awareness, IFUs, count sheets, status, location, and pickup \cite{steris-loaner-trays-2021}. & Loaner trays arrive with manufacturer-stocked contents and vendor documentation that uses the vendor's naming conventions and part numbers, not the hospital's local aliases or count-sheet format. SPD receives the tray with hours or minutes to process it before surgery. The worker must reconcile the vendor's nomenclature against the local count sheet or procedure card with no lead time to prepare or cross-reference \cite{steris-loaner-trays-2021}. \\
\hline
\end{tabularx}
\end{table}
\subsection{Environment}\label{sec:4-6-environment}
\begin{table}[H]
\centering
\caption{Fishbone: Environment root causes}
\label{tab:fishbone-environment-root-causes}
\small\renewcommand{\arraystretch}{1.2}\setlength{\tabcolsep}{4pt}
\begin{tabularx}{\textwidth}{|>{\raggedright\arraybackslash}X|>{\raggedright\arraybackslash}X|}
\hline
Observed failure mode & Root cause hypothesis \\
\hline
Alfred et al. identify production pressure and workspace constraints as performance-shaping factors in assembly work; AORN, Medline, and Mácola et al. describe short staffing, high turnover, rising demand, and financial stress as workforce pressures that reduce attention and training quality \cite{alfred-et-al-2021}; \cite{aorn-staffing-shortage-2025}; \cite{macola-et-al-2025}; \cite{brozak-2025}. Human factors such as stress and communication breakdowns between SPD and the OR compound the risk of instrument errors \cite{lind-2026a}. & SPD is a cost center within the hospital — it generates no direct revenue, so staffing budgets are structurally constrained relative to OR capacity (a revenue center). Hospitals routinely add surgery suites without proportionally increasing SPD headcount, as documented in a Colorado inspection report \cite{ingold-2025}. Rushed workers can't focus on tiny visual differences. Fewer staff means less time for training and less thorough checking \cite{alfred-et-al-2021}. \\
\hline
A Colorado hospital inspection report tied a major SPD backlog to increased staffing requirements after new operating rooms opened without matching SPD staff increases \cite{ingold-2025}. & When hospitals add surgery rooms without hiring more SPD staff, the work piles up. That backlog then makes every other problem — training pressure, checking shortcuts, rushing — even worse \cite{ingold-2025}. \\
\hline
Huang et al. examine training demands related to interruptions in central sterile supply departments \cite{huang-et-al-2025}. & Interruptions are structurally produced by the SPD workflow itself — phones ring with OR requests, surgeons call for add-on instruments, runners deliver and collect trays, delivery schedules interrupt assembly loops. The open-plan workstation layout makes it difficult to shield assembly workers from these interruptions \cite{huang-et-al-2025}. Getting interrupted while sorting or counting makes workers lose their place, skip items, and miss mistakes. \\
\hline
OSHA central sterile guidance discusses workstation, reaching, standing, carts, and height-adjustment concerns for central sterile work \cite{osha-central-sterile-supply-accessed-2026}. & Uncomfortable workstations, awkward reaches, and poor setup make a job that already needs close attention even harder \cite{osha-central-sterile-supply-accessed-2026}. \\
\hline
\end{tabularx}
\end{table}
\subsection{Policy And Feedback}\label{sec:4-7-policy-and-feedback}
\begin{table}[H]
\centering
\caption{Fishbone: Policy and feedback root causes}
\label{tab:fishbone-policy-and-feedback-root-causes}
\small\renewcommand{\arraystretch}{1.2}\setlength{\tabcolsep}{4pt}
\begin{tabularx}{\textwidth}{|>{\raggedright\arraybackslash}X|>{\raggedright\arraybackslash}X|}
\hline
Observed failure mode & Root cause hypothesis \\
\hline
HPN emphasizes standardized, clearly written count sheets that are available across shifts; Outpatient Surgery Magazine describes tray customization and standardization as reducing instrument identification time, and notes that tray information must be maintained as a shared local standard \cite{nadeau-2024}; \cite{loria-2024}. & Nobody owns the standards — SPD governance is decentralized across shifts, and no single role is responsible for maintaining, auditing, or enforcing count-sheet accuracy or tray-assembly procedure across all shifts. What you learn depends on who trains you that day, creating inconsistent mental models of the same tray among workers on the same line \cite{nadeau-2024}; \cite{loria-2024}. \\
\hline
Chobin states that instrument processing requires coordinated policies, procedures, accountability, education, and documentation \cite{chobin-2019}. & Without clear written rules for ambiguous situations — worn but functional instrument, substitute from a different manufacturer, unknown instrument in a loaner tray — workers rely on personal judgment, which varies by experience level and shift. No policy framework provides decision rules for these common edge cases \cite{chobin-2019}. \\
\hline
Sterile-processing certification commentary argues that certification creates a common baseline but cannot guarantee error-free work; Ofstead et al. used competency testing and booster training as a structured training intervention; managers must maintain competency programs and continuing education \cite{kovach-2012}; \cite{ofstead-et-al-2023}. Repeated errors in practice are often traceable to incomplete, outdated, or unused competency documents \cite{lind-2026b}. & If workers only get tested at hiring and never again, gaps in their knowledge, changes to trays, and worn instruments all stay hidden. Neither HSPA nor CBSPD mandates recurring competency checks tied to local tray knowledge — initial certification is not time-limited in a way that requires demonstrating current local competence \cite{kovach-2012}; \cite{ofstead-et-al-2023}. When competency programs exist, they are designed and administered locally with no external audit, so they range from rigorous to pro forma \cite{lind-2026b}. \\
\hline
Nichol et al. describe traditional surgical-instrument error reporting as cumbersome, human-dependent, delayed, incomplete, and difficult to integrate with operational systems; Zhu et al. observed packaging errors including wrong specifications, incomplete packages, and missing instruments — errors that could reflect upstream gaps not caught before packaging \cite{nichol-et-al-2024}; \cite{zhu-et-al-2019}. & When mistakes aren't caught and reported quickly, the same errors keep happening. The feedback loop from OR to SPD is broken — errors discovered during surgery are documented (if at all) in incident reports that may take weeks to reach SPD management. By that time the specific tray, technician shift, and training gap are hard to identify, so the error teaches no preventive lesson. Without closed-loop tracking that ties an error back to a root cause (individual, tray, training gap, or count-sheet error), downstream detection does not drive upstream improvement \cite{nichol-et-al-2024}; \cite{zhu-et-al-2019}. \\
\hline
\end{tabularx}
\end{table}

\section{Literature And Related Work}
\label{sec:5--literature-and-related-work}
The root cause analysis identified several interconnected causes of
tray-content errors. Among them, one knowledge-focused pattern recurs across
multiple failure modes: that general certification and broad instrument
familiarity may not reliably prepare technicians to recognize and differentiate
the full range of complex, lookalike, and locally variant instruments
encountered during tray reconstruction (Section~\ref{sec:4-2-product}). The same analysis also
identified product-level, process-level, information-level, and environmental
root causes, meaning there are many possible intervention points.
\subsection{Possible Solution Categories}
\label{sec:5-1-solution-landscape}
The solution space for reducing tray-content errors during reconstruction can
be organized into seven categories by intervention point. Each category targets
a different subset of the root causes, and each faces a
different feasibility constraint for this project.

The seven identified solution categories are:
\begin{enumerate}
  \item \textbf{Automated content verification}: post-assembly inspection using CV, RFID, or barcode scanning to verify tray contents.
  \item \textbf{Physical error-proofing}: modifying the tray or instrument interface (custom foam, color coding) to prevent selection errors.
  \item \textbf{Workflow / process design}: restructuring task handoffs and verification steps through dual checks or gating sign-offs.
  \item \textbf{Decision support at point of assembly}: providing standardized count sheets, photos, and alias lookup at the workstation.
  \item \textbf{Task simplification}: rationalizing tray contents and standardizing instrument sets across procedures.
  \item \textbf{Workforce / accountability systems}: competency reassessment, closed-loop error tracking, and structured mentorship.
  \item \textbf{Training and simulated practice}: building local tray skills through retrieval practice, simulation, and feedback before supervised work.
\end{enumerate}

\begin{table}[H]
\centering
\caption{Seven solution categories compared by constraint}
\label{tab:seven-solution-categories-compared-by-constraint}
\small\renewcommand{\arraystretch}{1.2}\setlength{\tabcolsep}{4pt}
\begin{tabularx}{\textwidth}{|>{\raggedright\arraybackslash}p{2.2cm}|>{\raggedright\arraybackslash}p{1.8cm}|>{\raggedright\arraybackslash}X|>{\raggedright\arraybackslash}p{3cm}|>{\raggedright\arraybackslash}X|}
\hline
Category & Intervention point & Example solutions & Root causes addressed & Key constraint for this project \\
\hline
Automated content verification & Post-assembly inspection & CV tray checker, RFID tray scan, barcode line scan, weight check & Product: wear hides differences and Process: no backup check exist & Requires real instruments, lighting, workflow integration; cannot validate without SPD deployment \\
\hline
Physical error-proofing & The tray or instrument interface & Custom foam cutouts, color-coded handles, segregated tray zones & Knowledge: workers must memorize tiny visual details and Process: assembly lacks a separation of pick from check & Expensive custom fabrication per tray; difficult to validate and maintain across tray variants \\
\hline
Workflow / process design & Task structure and handoffs & Dual verification, gating sign-offs, pre-sorted pick lists & Process: no structured backup check during assembly and Policy: nobody formally owns the standard & Primarily organizational change; no technical prototype to build or measure in a controlled study \\
\hline
Decision support at point of assembly & Information available during reconstruction & Standardized count sheets with photos, alias lookup, reference kiosk & Information: outdated count sheets force guessing and Knowledge: general instrument names do not convey local tray rules & Improves the reference but does not exercise or measure the technician's own ability \\
\hline
Task simplification & The tray definition itself & Tray rationalization, instrument consolidation, procedure-level standardization & Knowledge: differentiating lookalikes requires local familiarity and Environment: pressure and interruptions break focus & Requires hospital utilization data and multi-stakeholder buy-in; no buildable prototype \\
\hline
Workforce / accountability systems & Competency management and error feedback & Periodic competency reassessment, closed-loop error tracking, structured mentorship & Policy: standards are only tested at hiring and Feedback: slow error reporting lets mistakes repeat & Primarily organizational policy change; no measurable prototype within our project scope \\
\hline
Training and simulated practice & Pre-work and between-work skill building & VR simulation, digital flashcards, tangible tray simulation, structured OJT & Knowledge: lookalike discrimination requires repeated practice and Local knowledge: tray-specific names and rules are not taught in general curricula & Requires learning-science justification and a novice study, both accessible without SPD deployment \\
\hline
\end{tabularx}
\end{table}
Each category reduces a real risk. Categories 1 through 6, however, all depend
on resources that this project cannot access within our project timeline: SPD
deployment sites, real instrument inventories, hospital utilization data,
institutional policy authority, or multi-stakeholder buy-in. Category 7 —
training and simulated practice — is the only category whose primary evidence
can be collected with accessible participants (novices), low-cost materials
(printed cards), and a bounded pre/post study design on a simulated local tray
task. The remainder of this section examines the literature supporting each
relevant area in depth.
\subsection{SPD Error And Tray-Readiness Evidence}
\label{sec:5-2-spd-error-and-tray-readiness-evidence}
The error landscape described in Section~\ref{sec:3--problem-definition-and-scope} establishes the pattern: 55\% of tray defects occur during assembly (Alfred et al., 2021), visualization tasks account for 88.6\% of observed errors (Nichol et al., 2024), and packaging errors mask upstream problems (Zhu et al., 2019). TrayGuard targets the error categories that appear in that evidence: missing, wrong, extra, misidentified, wrong-count, and visibly unacceptable items. The current prototype focuses on the first five because they can be measured in a controlled card-based tray task.
\subsection{SPD Training And Competency Context}
\label{sec:5-3-spd-training-and-competency-context}
The RCA (Section~\ref{sec:4-3-knowledge}) identified the gap between general certification and local tray competence: HSPA's CRCST requires 400 hours of hands-on experience, but certification alone does not cover local instrument names, tray rules, or sustained lookalike discrimination \cite{hspa-crcst-accessed-2026}. CBSPD's CSPDT route similarly treats certification as minimum competency, allowing candidates to qualify through work experience, allied health plus SPD experience, a sterile processing course, or related product sales or service experience \cite{cbspd-technician-accessed-2026}.

SPD training in practice follows a blended pattern that combines coursework, hands-on practice, and preceptor-led onboarding:
\begin{table}[H]
\centering
\caption{SPD training patterns and TrayGuard implications}
\label{tab:spd-training-patterns-and-trayguard-implications}
\small\renewcommand{\arraystretch}{1.2}\setlength{\tabcolsep}{4pt}
\begin{tabularx}{\textwidth}{|>{\raggedright\arraybackslash}p{3cm}|>{\raggedright\arraybackslash}X|>{\raggedright\arraybackslash}X|}
\hline
Training pattern & Evidence & Implication for TrayGuard \\
\hline
Certification exam plus verified hands-on hours & HSPA CRCST requires 400 SPD hands-on hours and a computer-based exam; CBSPD CSPDT uses eligibility routes that include work experience or a course plus exam \cite{hspa-crcst-accessed-2026}; \cite{cbspd-technician-accessed-2026}. & The tool should measure learning and support practice before and during supervised competency development. \\
\hline
Self-paced or classroom instruction & Towson's ed2go-backed course is self-paced online, 190 hours, covering instrument identification and tray assembly. Altamont's program is 135 hours over 8-12 weeks with live online classes, in-person labs, and an optional 240-400 hour externship \cite{towson-sterile-processing-accessed-2026}; \cite{altamont-healthcare-accessed-2026}. & The platform can complement formal instruction with reusable local modules and repeated practice. \\
\hline
Preceptor-led onboarding & HPN describes hospital examples including 90-day on-the-job training programs, 19-week programs with textbook chapters, quizzes, and progress tests, and structured 1:1 preceptor onboarding with presentations, online modules, written tests, direct observation, and one-on-one training \cite{nadeau-2017}. HSPA training guidance frames core program elements as structured orientation, competency documentation, and ongoing education; training that stops after orientation perpetuates inconsistent habits \cite{klacik-2023}. & The system should reduce low-level explanation burden and give preceptors evidence of where a learner still struggles. \\
\hline
Simulation and skills assessment & AORN describes Penn Medicine's earn-to-learn SPD program combining didactic education, hands-on training, simulation, mentorship, content testing, and skills-based assessment \cite{aorn-staffing-shortage-2025}. & Simulated tray sorting and pre/post tasks are aligned with existing training practice. \\
\hline
\end{tabularx}
\end{table}

The common thread across patterns is that hands-on practice and preceptor attention are the scarce resources. Several program examples show the gap between classroom preparation and clinical capacity. Penn Foster's online sterile processing program prepares learners for the CRCST exam but does not require the exam or the 400 hands-on hours \cite{penn-foster-sterile-processing-accessed-2026}. MedCerts states that full CRCST certification requires 400 hours of hands-on experience following provisional certification \cite{medcerts-catalog-2025}. AIMS Education requires a 400-hour clinical internship, but its schedules depend on clinical-site availability \cite{aims-education-sterile-processing-accessed-2026}. Central Sterilization Solutions states its classroom course has no hands-on training in the classroom and that hands-on training is conducted during externship \cite{central-sterilization-solutions-in-person-accessed-2026}. These examples illustrate the training-access gap: repeatable, feedback-rich practice with local instruments and tray rules is hard to scale before and during real SPD experience because it depends on site access, clinical scheduling, preceptor attention, and available teaching materials.

Broader allied-health placement research supports treating clinical placements as capacity-constrained environments. A Queensland study found that placement-offer growth did not keep pace with student-program growth, and barriers centered on resourcing and increasing demand \cite{mcbride-et-al-2020}. A 2025 Greater Sacramento report identified shortage of clinical placement opportunities as a barrier to workforce pipeline growth, and highlighted student preparedness as part of successful placement partnerships \cite{north-greater-sacramento-coe-2025}. TrayGuard gives learners repeated practice before, during, or between supervised work, and gives educators evidence about weak instruments, weak tray rules, high-confidence errors, and uncertainty. This addresses the practical constraint: supervised training time is a scarce resource \cite{aorn-staffing-shortage-2025}; \cite{brozak-2025}; \cite{macola-et-al-2025}.
\subsection{Local Tray And Count-Sheet Evidence}
\label{sec:5-4-local-tray-and-count-sheet-evidence}
The RCA (Section~\ref{sec:4-5-information}) identified incomplete count sheets and local tray variation as root causes. Count sheets define tray names, sections, contents, quantities, reference numbers, placement, and sign-off fields (Nadeau, 2024; Stryker, 2023). Tray contents are design choices shaped by procedure, surgeon, and hospital policy (dos Santos et al., 2021; Ahmadi et al., 2023), so general instrument knowledge does not transfer across sites. Loaner trays add further complexity through scheduling, delivery, IFU, and count-sheet dependencies (STERIS, 2021). TrayGuard stores local names, aliases, approved photos, required counts, lookalike pairs, distractors, assessment variants, and module version. Modules must be instructor-verified before results are described as content-valid.

FGVC-derived modules serve a narrower prototype role. Atabuzzaman et al. collected ultra-fine-grained surgical instrument images from major, minor, and eye tray contexts \cite{atabuzzaman-et-al-2025}. TrayGuard's current FGVC prototype uses the major-tray image archive and class-name mapping to create image cards and same-family discrimination exercises. The FGVC source provides observed instrument classes and representative images; public count sheets provide the field structure for tray names, sections, quantities, locations, check boxes, and sign-off fields \cite{stryker-gamma4-count-sheet-2023}; \cite{stryker-imn-count-sheet-2023}. The prototype's FGVC quantities are simulated as one card per selected class, and exact local tray membership remains an instructor-review variable.
\subsection{Learning Science And Simulation Evidence}
\label{sec:5-5-learning-science-and-simulation-evidence}
Health professions education provides a general precedent for simulation: a
broad review of technology-enhanced simulation found improved knowledge, skills,
and behaviors compared with no intervention, and simulation-based medical
education with deliberate practice has also been shown to outperform traditional
clinical education in meta-analysis
\cite{cook-et-al-2011};
\cite{mcgaghie-et-al-2011}. Earlier medical
simulation guidance identifies feedback, repetitive practice, curriculum
integration, and measurable outcomes as important features of effective
simulation \cite{issenberg-et-al-2005}.
For sterile processing specifically, Ofstead et al. provide the closest direct
precedent. Their borescope/endoscope visual-inspection training pilot used a
pre-test, structured teaching, hands-on practice, image-based testing,
confidence/satisfaction measures, workplace homework, and a delayed booster
with certified sterile-processing professionals
\cite{ofstead-et-al-2023}. Tray assembly is
a different task, but the learning structure maps well: the learner must build
visual familiarity, recognize subtle differences, decide when uncertain, apply
a procedural rule, and improve between a baseline task and a later task.
Retrieval-practice evidence explains why the prototype should ask learners to
produce answers before showing them. Tests can improve long-term retention
relative to restudy, repeated testing improved long-term retention in a
medical-education randomized trial, and reviews support practice testing,
distributed practice, spaced digital education, electronic flashcards, and
actionable feedback in health-professions learning
\cite{roediger-and-karpicke-2006};
\cite{larsen-et-al-2009};
\cite{dunlosky-et-al-2013};
\cite{martinengo-et-al-2024};
\cite{barrison-et-al-2025};
\cite{hattie-and-timperley-2007}.
This supports retrieval-first cards, quizzes, weak-item repetition, and
error-specific feedback as the instructional core rather than decoration around
the tray task.
The electronic-flashcard evidence is especially relevant to the proposed study
and card modes. Barrison et al.'s 2025 scoping review found that electronic
flashcards are widely used in health-professions education, that the research
base grew rapidly from 2019 to 2024, and that many studies examine utilization
and associated learning outcomes. The review also cautions that development and
delivery methods are less systematically studied
\cite{barrison-et-al-2025}. For TrayGuard,
that means flashcards are a credible medical-learning mechanism, but they
should be designed as one part of a larger tray simulation: short prompts,
single target concepts, local photos or card representations, immediate answer
feedback, spacing or weak-item repetition, and a downstream tray-sort task that
tests application rather than isolated recall.
The proposed AprilTag-card tray is a reasonable physical proxy because it
preserves the target cognitive actions while replacing regulated instruments
with safe, cheap, repeatable tokens. The user still has to read a local tray
list, retrieve which items belong, distinguish distractors and lookalikes, apply
counts, place items into a tray area, and respond to feedback. The marker
makes the simulated choice observable to the app. NeuroVase is a close emerging
design precedent: it uses tangible cue cards with a tablet-based mobile AR
system, structured medical curriculum, pre/post assessment, usability measures,
and a controlled user study for neurovascular anatomy and stroke education
\cite{jahani-et-al-2026-neurovase}.
QR-code education literature similarly supports low-cost printable codes as an
access layer for healthcare learning, simulation, and training support
\cite{karia-et-al-2019}.
The same design pattern appears in workplace training research.
Meta-analytic evidence shows that organizational training improves learning,
behavior, and operational results when it includes practice and feedback
\cite{arthur-et-al-2003}, and Salas et al. argue that effective training should
be treated as a system: analyze the work, design instruction around target
competencies, provide practice and feedback, and evaluate transfer
\cite{salas-et-al-2012}. TrayGuard follows that model by targeting a specific
work task (tray reconstruction), designing instruction around instrument
identity and count-sheet rules, providing repeated practice with error-specific
feedback, and measuring pre/post transfer to a comparable tray variant.
Serious tabletop and board-game precedents also support the choice to turn a
healthcare protocol into a bounded physical learning game. Ward et al. designed
and evaluated the PlayDecide patient-safety board game in two acute teaching
hospitals; the intervention used cards and facilitated discussion to teach
junior doctors about safety concerns and reporting, and the authors concluded
that it was valuable for patient-safety education and deep discussion
\cite{ward-et-al-2019}. Their study demonstrates
that abstracted physical artifacts can stand in for clinical objects when the
study measures learning, discussion, decision quality, and protocol use.
Marker cards are the first prototype for measuring novice improvement on
simulated local tray sorting. A follow-on SPD pilot would replace or supplement
cards with local photos, real teaching instruments, educator-reviewed count
sheets, and supervised workplace validation.
\subsection{Computer Vision Feasibility And Boundary Evidence}
\label{sec:5-6-computer-vision-feasibility-and-boundary-evidence}
Computer vision remains technically plausible, but it is no longer the central
project proof. Deol et al. show that automated surgical-instrument detection and
counting can work in an experimental proof-of-concept setting
\cite{deol-et-al-2024}. Atabuzzaman et al.
show that structured image acquisition can support ultra-fine-grained surgical
instrument classification, and Xin et al. address CSSD-oriented instrument
recognition
\cite{atabuzzaman-et-al-2025};
\cite{xin-et-al-2024}. These papers support CV as authoring support, camera-assisted practice, or a
future visual review layer, but real SPD deployment would need local data,
manufacturer and variant coverage, lighting and glare testing, occlusion
handling, open-set behavior, workflow review, human confirmation, false-alert
analysis, cleaning and device-policy review, and site-specific failure-mode
reporting. Kienle et al. show why this boundary
matters: strong in-domain detection can drop when instruments come from
different manufacturers
\cite{kienle-et-al-2025}.
The project's earlier CV pipeline is still useful technical evidence, but it
should be treated as a support layer. The project's primary claim should come from
the training loop because that loop can be built, tested, and interpreted
without pretending that deployment-grade instrument recognition has already
been solved.
\subsection{Market And Comparable Tools}
\label{sec:5-7-market-and-comparable-tools}
Comparable systems cluster into four groups: tray-management products,
instrument tracking systems, surgical simulation tools, and emerging CV or
robotic research systems. None of them eliminates the need for TrayGuard's
bounded training claim.
\begin{table}[H]
\centering
\caption{Comparable tools and gap relative to TrayGuard}
\label{tab:comparable-tools-and-gap-relative-to-trayguard}
\small\renewcommand{\arraystretch}{1.2}\setlength{\tabcolsep}{3pt}
\begin{tabularx}{\textwidth}{|>{\raggedright\arraybackslash}p{2.5cm}|>{\raggedright\arraybackslash}p{2.5cm}|>{\raggedright\arraybackslash}X|>{\raggedright\arraybackslash}p{2cm}|>{\raggedright\arraybackslash}X|}
\hline
Tool or approach & Target user & Relevant features & Evidence strength & Gap relative to TrayGuard \\
\hline
Tray Pacer & SPD teams and managers & Count sheets, tray assembly support, photos, productivity and proficiency metrics. & Commercial product evidence. & Supports operations, but public materials do not establish the specific pre/post local tray learning study proposed here \cite{tray-pacer-accessed-2026}. \\
\hline
LayerJot SID & SPD and surgical instrument teams & Instrument documentation, photos, count sheets, and tray information. & Commercial product evidence. & Strong local-information precedent, but not a demonstrated tag-card novice learning loop \cite{layerjot-sid-accessed-2026}. \\
\hline
CensiTrac and similar tracking systems & Hospitals, SPD, OR inventory teams & Instrument and tray tracking, traceability, inventory visibility. & Commercial and operational precedent. & Addresses tracking and lifecycle visibility more than low-cost novice simulation and weak-item training \cite{censitrac-accessed-2026}. \\
\hline
Touch Surgery and surgical simulators & Medical trainees and clinicians & Repeated simulation attempts, performance logging, skill practice. & Published evaluation and commercial adoption. & Supports simulation logic, but focuses on surgical procedure skills rather than SPD tray reconstruction \cite{tulipan-et-al-2019}; \cite{medtronic-touch-surgery-accessed-2026}. \\
\hline
NeuroVase-style tangible cards & Health-professions learners & Physical cue cards, mobile AR, structured curriculum, pre/post testing, usability measures. & Emerging preprint evidence. & Strong design precedent for tangible medical learning, but not sterile-processing content \cite{jahani-et-al-2026-neurovase}. \\
\hline
Surgical-instrument CV research & Researchers and future technical teams & Detection, counting, fine-grained classification, instrument-stand recognition. & Published research, usually bounded datasets. & Supports future camera assistance, but not deployment-ready SPD workflow claims \cite{deol-et-al-2024}; \cite{atabuzzaman-et-al-2025}; \cite{kienle-et-al-2025}. \\
\hline
Autonomous tray assembly research & Robotics researchers and future automation teams & Manipulation, sorting, and assembly concepts. & Early research. & Much higher hardware, sterility, safety, and integration burden than the project's training loop \cite{da-silva-et-al-2026}. \\
\hline
\end{tabularx}
\end{table}
TrayGuard's contribution is a low-cost local training loop that turns a tray
module into retrieval-first study, applied tray sorting, immediate feedback,
no-hints assessment, and exportable pre/post evidence.

Medical training software already uses the same design pattern of simulation,
guided practice, feedback, scoring, and progress tracking. Touch Surgery
provides mobile surgical procedure simulations with repeated attempts and
scoring, and a validation study found that users improved across repeated
attempts while higher-experience cohorts outperformed lower-experience cohorts
\cite{tulipan-et-al-2019}. Osso VR demonstrated in a randomized trial that
VR-trained novices completed a downstream orthopedic procedure more often, made
fewer incorrect steps, and finished faster than guide-only learners
\cite{orland-et-al-2020}. Body Interact uses screen-based virtual patient
scenarios with immediate feedback, auto-grading, and educator analytics
\cite{wolters-kluwer-body-interact-accessed-2026}. Surgical Science and
Simbionix provide commercial medical simulators across specialties,
positioning simulation-based assessment as a credible healthcare training
category \cite{surgical-science-simulators-accessed-2026}. These precedents
support treating TrayGuard as a health-professions simulation tool that applies
established learning mechanisms to a new domain: local SPD tray reconstruction.
\subsection{Literature Gaps and System Response}
\label{sec:5-8-literature-gaps-and-system-response}
The literature supports each design choice individually, but no existing work
closes the specific gap that TrayGuard targets: a low-cost, repeatable,
locally configurable training loop for SPD tray reconstruction that measures
pre/post simulated improvement.
\begin{table}[H]
\centering
\caption{Literature gaps and TrayGuard's system response}
\label{tab:literature-gaps-and-system-response}
\small\renewcommand{\arraystretch}{1.2}\setlength{\tabcolsep}{4pt}
\begin{tabularx}{\textwidth}{|>{\raggedright\arraybackslash}X|>{\raggedright\arraybackslash}X|>{\raggedright\arraybackslash}X|}
\hline
Gap in the literature & System response & Stakeholder meaning \\
\hline
CV papers show instrument-detection feasibility while leaving manufacturer, site, tray, and workflow variation thinly tested \cite{deol-et-al-2024}; \cite{atabuzzaman-et-al-2025}; \cite{kienle-et-al-2025}. & Keep CV optional and locally verifiable. Use instructor-approved content as the source of truth for the training evidence. & The tool does not depend on deployment-grade CV. Camera support can be added later after site-specific validation. \\
\hline
SPD error papers document real defect rates while leaving the training-intervention mechanism unresolved \cite{alfred-et-al-2021}; \cite{nichol-et-al-2024}. & Run a pre/post simulated learning study with error-category breakdown. & Students can test novice learnability; SPD trainees remain the future validation group for workplace transfer. \\
\hline
Training literature supports retrieval practice, simulation, spacing, and feedback while leaving the SPD tray-module design open \cite{roediger-and-karpicke-2006}; \cite{cook-et-al-2011}; \cite{ofstead-et-al-2023}. & Build the module around local retrieval cards, quizzes, applied sorting, weak-item repetition, and error-specific feedback. & Trainees get repeatable practice; instructors get evidence about weak items. \\
\hline
Adoption literature provides broad healthcare-AI guidance without SPD training-assistant specifics \cite{goddard-et-al-2012}; \cite{kelly-2026}. & Report workflow friction, role boundaries, privacy assumptions, and pilot requirements alongside the prototype. & Administrators get a clearer approval path and an explicit future-evidence list. \\
\hline
\end{tabularx}
\end{table}
\subsection{Final Product Concept}
\label{sec:6-1-final-product-concept}
TrayGuard is a local tray training and assessment platform. The core product is
a learner-facing app plus a file-backed tray module that defines the
instruments, aliases, required counts, distractors, lookalike pairs,
distinguishing features, and assessment variants for one local tray.
The first build should use printable AprilTag instrument cards as the physical
simulation layer. Each card represents one instrument or instrument variant.
The app scans the cards on a tabletop or tray area, treats the detected markers
as the learner's tray selection, scores the result against the tray template,
and returns missing, extra, wrong, misidentified, and wrong-count feedback in
practice mode. Pre-test and post-test modes suppress hints and corrective
feedback so the same system can collect baseline and post-training evidence.
The full learning loop is:
\begin{lstlisting}
local tray module -> pre-test -> retrieval-first cards -> quiz -> AprilTag-card tray sort -> post-test -> metrics export
\end{lstlisting}
This concept keeps the prototype buildable while preserving the important task
structure: the learner must use local names, local counts, visual/semantic
distinctions, and tray rules to assemble a simulated tray.
\section{Final Product Concept And Design Rationale}
\label{sec:6--final-product-concept-and-design-rationale}
The final TrayGuard direction is a local tray training and assessment platform
that addresses tray reconstruction errors by reducing training and
knowledge-transfer risk. The selected concept moves the proof point earlier in
the workflow: can a novice learn a local tray module well enough to improve on
a comparable simulated tray-sorting task?
\subsection{Rationale For The Training-Centered Pivot}
\label{sec:6-2-rationale-for-the-training-centered-pivot}
The project originally aimed at a live CV tray checker: detect instruments,
compare detections against a tray list, and return a readiness result. That
direction remains attractive, but it has the wrong proof burden for our current
project. A useful clinical tray checker would need strong evidence across real
instruments, local tray definitions, similar variants, glare, occlusion,
unknown objects, workflow interruptions, false alerts, human review, device
policy, cleaning, privacy, and adoption. The current project can demonstrate
parts of that pipeline, but it cannot responsibly claim deployment-grade
clinical reliability.
The Week 6 pivot made the value proposition more testable: reduce simulated
time-to-competency and training burden for sterile-processing learners by
building an interactive training platform. That direction keeps the same root
problem, because missing, wrong, extra, and misidentified instruments still
arise during tray reconstruction. It changes the claim from "the model can
outperform humans in a hospital environment" to "novices can improve on a
measured local tray simulation after retrieval-centered practice."
This is a better project claim for three reasons. First, it matches the evidence:
sterile-processing work depends on local knowledge and supervised practice, and
simulation/retrieval evidence supports repeated practice with feedback. Second,
it is buildable with low-cost materials: a file-backed tray module, printable
cards, a scanning surface, scoring, and export. Third, it reduces risk for a
future CV system by creating validated local content, error categories,
workflow language, and learner data before asking a camera model to make
clinical-grade distinctions.
\subsubsection{What The CV Experiments Showed}
\label{sec:6-2-1-what-the-cv-experiments-showed}
The CV experiments were structured around three failure dimensions identified in
the literature: lighting sensitivity, lookalike confusion, and domain transfer.
Each dimension was tested with YOLO variants (YOLO11n/s/m/l) on self-collected
and synthetic datasets before the pivot decision.
\paragraph{Lighting robustness.}
Seven staged experiments measured how detection degrades when training and test
illumination differ. Models trained at 800 lumens and tested at 80 lumens showed
monotonic mAP50-95 decay as brightness decreased. Models trained on dark images
and tested on bright images showed a similar but asymmetric decay. Background
transfer (matte-to-reflective, reflective-to-matte) caused additional mAP50-95
loss at every brightness level. A separated-to-overlay layout shift, where
training images showed instruments in isolation and test images showed
instruments overlapping, produced the largest drop: mAP50-95 fell below 0.3
across all brightness ranks. The gap between the separated baseline and the
overlay-test condition was consistent enough to conclude that training on
neatly arranged instruments does not transfer to the cluttered arrangements
that occur when multiple instruments occupy a single tray view.
These results showed that controlled lighting and layout are prerequisites for
reliable detection, and that real SPD workstations produce variable lighting,
reflective surfaces, and overlapping instruments that would require
site-specific training data to cover.
\paragraph{Shape similarity and pairwise confusion.}
Six instrument-like classes organized into three pair families (straight versus
curved scissors, straight versus curved clamps, narrow-tip versus broad-tip
tweezer/plier) were rendered in Blender and photographed as real proxies.
Pairwise confusion analysis showed that same-family class pairs were the most
common error mode: straight scissors were regularly detected as curved scissors
and vice versa, with confusion rates exceeding 30\% of predictions in the
synthetic-seen-condition evaluation. High-confidence wrong predictions
(confidence $\ge$ 0.60) accounted for a measurable fraction of these errors,
meaning the model was confidently wrong on fine-grained distinctions that are
exactly the kind of lookalike pairs that SPD technicians must learn to separate.
Open-set evaluation added unknown-similar-distractor objects to the test split.
The unknown false positive rate, the fraction of unknown images that triggered
at least one spurious detection, reached 40--60\% across both the synthetic
and real evaluation stages. The model did not reject unfamiliar objects; it
assigned them a nearest known class label with moderate to high confidence.
This directly undermines the value of a deployment tray checker, because
unexpected instruments, vendor loaner items, or worn variants would produce
false detections rather than explicit uncertainty.
\paragraph{Synthetic-to-real transfer.}
A small real-proxy dataset (roughly 200 images across the six classes) was used
to compare training from scratch against fine-tuning from a synthetic pretrained
checkpoint. The synthetic pretraining improved mAP50 by approximately 5--10
percentage points over training from scratch, but the absolute mAP50 on the
real test set remained below 0.70. Pairwise confusion and unknown false positive
rates did not improve proportionally. The pretraining benefit was real but
bounded: synthetic data helped initialize feature representations but did not
solve the fine-grained discrimination or open-set rejection problems that a
deployment checker requires.
\paragraph{What the cumulative evidence meant for the project.}
Eight trained weight files across YOLO11n/s/m/l variants, RT-DETR-l, and
ablated lighting conditions were evaluated on multiple dataset stages. The
pattern across all experiments was consistent: detection worked in controlled,
matched-condition settings and degraded under every variation that SPD
workflows routinely present -- different lighting, reflective backgrounds,
overlapping instruments, lookalike variants, and unknown items. The project
could have continued collecting more data, training more models, and narrowing
the operating conditions, but each tightened constraint would have moved the
project further from the stated goal of a practical tray-verification tool.
The pivot to a training-centered platform preserved the technical infrastructure
and reused the data collection pipeline, camera code, and detection experiments
as foundation for future camera-assisted practice. The primary evidence path
shifted from CV generalization to measurable novice learning through a
simulated tray task with physical AprilTag cards, a claim that could be
supported with the project's available resources.
\subsection{Alternatives Considered}
\label{sec:6-3-alternatives-considered}
Several product directions were considered:
\begin{table}[H]
\centering
\caption{Alternative product directions considered and rejected}
\label{tab:alternative-product-directions-considered-and-rejected}
\small\renewcommand{\arraystretch}{1.2}\setlength{\tabcolsep}{4pt}
\begin{tabularx}{\textwidth}{|>{\raggedright\arraybackslash}X|>{\raggedright\arraybackslash}X|>{\raggedright\arraybackslash}X|}
\hline
Alternative & Benefit & Reason rejected or deferred \\
\hline
Deployment-grade CV tray checker on real instruments & Closest to the original automation concept and strongest if validated in an SPD. & Requires local instrument data, clinical workflow integration, failure-mode validation, and regulatory/adoption work beyond our project timeline. \\
\hline
Software-only study and quiz module & Fastest implementation path and easiest novice study. & Too far from the tray reconstruction task; it would mostly test recognition, not applied count-sheet use. \\
\hline
Real teaching instruments with camera-supported tray review & Better physical fidelity than cards. & Requires access to representative instruments, cleaning/handling rules, storage, and educator review; appropriate for an SPD pilot after the simulated loop works. \\
\hline
Printable QR/AprilTag instrument cards & Cheap, portable, repeatable, and compatible with device scanning; closely matches NeuroVase-style tangible cue-card learning. & Less physically realistic than real instruments, so claims must stay limited to simulated tray-task improvement. \\
\hline
Manual drag-and-drop tray simulation & Simple software implementation without camera setup. & Does not exercise physical search, placement, or scanning workflow, and is less persuasive for an interactive project demo. \\
\hline
\end{tabularx}
\end{table}
The selected first prototype is the printable AprilTag-card path because it
creates a physical, measurable simulation without depending on deployment-grade
instrument recognition.
\subsection{Selected Candidate And Backup Paths}
\label{sec:6-4-selected-candidate-and-backup-paths}
The selected candidate is:
\begin{quote}
A local tray training and assessment app that uses printable AprilTag
instrument cards for simulated tray sorting, retrieval-first study, quiz,
practice feedback, no-hints pre/post assessment, and metrics export.
\end{quote}
The backup paths are layered so the project still produces evidence if one
piece fails.
\begin{table}[H]
\centering
\caption{Backup path strategy by risk level}
\label{tab:backup-path-strategy-by-risk-level}
\small\renewcommand{\arraystretch}{1.2}\setlength{\tabcolsep}{4pt}
\begin{tabularx}{\textwidth}{|>{\raggedright\arraybackslash}X|>{\raggedright\arraybackslash}X|>{\raggedright\arraybackslash}X|>{\raggedright\arraybackslash}X|}
\hline
Level & Selected path & Backup path & Claim preserved \\
\hline
Learning feature & Retrieval cards, quiz, practice sort, pre/post sort. & If the full loop is too long, run pre-test, study cards, practice sort, and post-test only. & Immediate simulated learning. \\
\hline
Physical interface & AprilTag cards scanned on a tray mat. & Use QR codes or manual card selection if AprilTag scanning is unreliable. & Local tray-task learning, with marker reliability reported as a limitation. \\
\hline
Content & \texttt{basic\_general\_tray\_v1} with 8-12 required concepts, distractors, lookalikes, and matched variants. & Reduce to 6-8 required concepts while preserving at least two lookalike pairs and count errors. & Bounded module feasibility. \\
\hline
CV & Marker detection only for project scoring. & Defer all camera recognition and score user selections manually. & Training evidence without CV overclaim. \\
\hline
Evaluation & 8-12 novice participants plus delayed retention if schedule allows. & Run a smaller formative walkthrough and report it as usability/protocol evidence, not learning proof. & Risk-reduction handoff. \\
\hline
Content validity & SPD educator or instructor review. & Faculty or surgical-technology reviewer plus explicit future SPD review requirement. & Plausibility review, not clinical validation. \\
\hline
\end{tabularx}
\end{table}
This backup strategy is intentionally conservative. It protects the project from
collapsing into a fragile detector demo and keeps the main deliverable aligned
with the training pivot.
\subsection{Rationale for AprilTag Marker Choice}
\label{sec:6-5-rationale-for-apriltag-marker-choice}
The prototype uses printable AprilTag fiducial markers on physical instrument
cards. This section justifies AprilTag specifically against the alternatives of
no physical component, QR codes, ArUco markers, and real instruments with
camera-based recognition.
\subsubsection{Why Physical Markers at All}
\label{sec:6-5-1-why-physical-markers-at-all}
A purely digital drag-and-drop tray simulation would be simpler to implement
and would avoid camera-setup and detection-reliability concerns. It is rejected
because the target task is fundamentally physical: SPD technicians search for,
pick up, examine, count, and arrange real instruments under a tray boundary
\cite{nichol-et-al-2024};
\cite{alfred-et-al-2021}. A purely on-screen
interface removes the spatial search, the need to distinguish similar objects
in hand, and the physical count-and-place workflow that the study is meant to
simulate. Physical cards preserve those cognitive demands while replacing
regulated stainless-steel instruments with safe, cheap, tractable tokens.
\subsubsection{Why AprilTag over QR Codes}
\label{sec:6-5-2-why-apriltag-over-qr-codes}
QR codes are a natural alternative because they are widely available, readable
by consumer devices, and familiar to most users. However, they are designed for
data encoding rather than computer-vision pose estimation and marker tracking
\cite{opencv-aruco-detection-accessed-2026}. A QR
code embeds dozens to hundreds of bytes of payload data, decodes it through a
Reed-Solomon error-correction pipeline, and returns a string. AprilTag inverts
that priority: it encodes a compact integer identifier using a binary matrix
optimized for fast corner detection, rotation disambiguation, and error
identification under perspective distortion, partial occlusion, and varying
lighting \cite{olson-2011}.
Three practical differences make AprilTag the better choice for this prototype.
First, detection speed and robustness. AprilTag exploits its known black border
and grid structure to extract corners at sub-pixel precision. The OpenCV
\texttt{aruco} module (which natively supports AprilTag dictionaries) can detect and
decode dozens of markers in a single video frame in milliseconds, while QR-code
detection requires finding three finder patterns, performing perspective
correction, and decoding a payload — a heavier pipeline designed for different
goals \cite{garrido-jurado-et-al-2014};
\cite{romero-ramirez-et-al-2018}.
Experimental comparisons have shown that AprilTag achieves higher detection
rates than QR codes under motion blur, shadow, and perspective distortion in
pose-estimation contexts
\cite{kalaitzakis-et-al-2021}.
Second, compact dictionaries with guaranteed Hamming-distance bounds. AprilTag
uses a lexicode-based coding system that guarantees a minimum Hamming distance
between every pair of markers under all four rotations
\cite{olson-2011};
\cite{wang-et-al-2016-apriltag}. A high
inter-marker distance means the detector can confidently identify a marker even
when several bits are corrupted by glare, print quality, or partial occlusion.
For a training prototype that must not confuse a learner's correct selection
with a detection error, that robustness is critical.
Third, pose estimation. A single AprilTag marker provides four coplanar corner
correspondences, sufficient for full 6-DOF pose estimation from a single
calibrated camera \cite{olson-2011}. Although the
current prototype does not yet use pose for detailed spatial feedback, the same
markers can later support richer interaction — detecting which instrument a
learner hesitates over, flagging instruments placed in the wrong tray zone, or
verifying the physical arrangement against a placement template. QR codes do
not offer comparable corner-localization accuracy because their finder patterns
are designed for barcode scanning rather than metric pose estimation.
\begin{table}[H]
\centering
\caption{AprilTag versus QR code comparison}
\label{tab:apriltag-versus-qr-code-comparison}
\small\renewcommand{\arraystretch}{1.2}\setlength{\tabcolsep}{4pt}
\begin{tabularx}{\textwidth}{|>{\raggedright\arraybackslash}X|>{\raggedright\arraybackslash}X|>{\raggedright\arraybackslash}X|>{\raggedright\arraybackslash}X|}
\hline
Requirement & QR code & AprilTag & Effect on TrayGuard \\
\hline
Detection under perspective, partial occlusion, and lighting variation & Moderate; designed for fronto-parallel scanning & Strong; designed for arbitrary camera pose and varying conditions & Fewer detection failures that inflate learner error rates \cite{kalaitzakis-et-al-2021}. \\
\hline
Sub-pixel corner accuracy & Not designed for pose estimation & Native: lexicode-guaranteed dictionary with sub-pixel refinement & Enables future tray-zone and placement feedback \cite{romero-ramirez-et-al-2018}. \\
\hline
Compact identifier with error correction & Large payload, Reed-Solomon on data bytes & Compact ID with lexicode Hamming-distance guarantee & Lower decode latency; no ambiguity between marker ID and payload parsing \cite{olson-2011}. \\
\hline
Open-source library in existing stack & OpenCV and platform libraries & Integrated in OpenCV \texttt{aruco} module; standalone C library & Zero additional dependency cost \cite{opencv-aruco-detection-accessed-2026}. \\
\hline
Neutral card front for assessment & QR pattern is visually distinctive and recognizable as a code & Binary grid is less obviously a machine-readable code to novices & Lower risk of learners recognizing markers as answer keys during assessment. \\
\hline
\end{tabularx}
\end{table}
\subsubsection{Why AprilTag over ArUco}
\label{sec:6-5-3-why-apriltag-over-aruco}
ArUco is the closest alternative: both produce binary square fiducial markers,
both are open source, and both can be detected through the same OpenCV \texttt{aruco}
module. For this project's requirements — up to 1000 markers and high detection
accuracy — AprilTag is the stronger choice for three reasons.
First, AprilTag's lexicode system provides mathematical guarantees that ArUco's
search-based dictionary does not. The lexicode algorithm generates a family of
codewords with a proven minimum Hamming distance under every possible rotation
\cite{olson-2011};
\cite{wang-et-al-2016-apriltag}. This means the
detector can not only detect but also correct more bit errors per marker than
ArUco's heuristic dictionary of equivalent size. A larger minimum distance
directly improves detection accuracy in the conditions most likely to occur
during a training session: glare off laminated cards, low tabletop lighting,
partial card occlusion, and off-angle camera views.
Second, AprilTag's corner-refinement pipeline (especially AprilTag 2's
adaptive thresholding and edge refinement) produces more stable corner
localization across varying lighting and image scales
\cite{wang-et-al-2016-apriltag}. The original
AprilTag was already benchmarked as the most accurate fiducial system in
head-to-head comparisons under motion blur and perspective distortion
\cite{kalaitzakis-et-al-2021}.
Third, AprilTag 3 introduces a multi-threaded detector and flexible tag layouts
that improve throughput on higher-resolution images without sacrificing
detection rate
\cite{krogius-et-al-2019}. For the first
build, the recommended family is \texttt{tagStandard52h13}, which provides 52 data
bits and a minimum Hamming distance of 13 — enough error-correction capacity
to support thousands of unique markers while still fitting comfortably on a
printed card. If fewer than 500 markers are needed, \texttt{tagStandard41h12} is a
more compact alternative with comparable robustness. Both families are
supported through the AprilTag dictionary support in OpenCV's \texttt{aruco} module,
with full backward compatibility for earlier dictionaries
(\texttt{DICT\_APRILTAG\_36h11}, \texttt{DICT\_APRILTAG\_16h5}), so no additional library
dependency is required.
\subsubsection{Why Not Real Instruments}
\label{sec:6-5-4-why-not-real-instruments}
The most obvious alternative — and the one that keeps the project closest to its
original CV-checker goal — is to use real surgical instruments with a
camera-based object detector. That path is rejected for the first prototype for
four reasons. First, real instruments require regulated access, cleaning and
handling protocols, storage space, and educator supervision that the first
study cannot assume \cite{hspa-crcst-accessed-2026}.
Second, instrument recognition across manufacturers, wear states, and lighting
conditions is an unsolved research problem that would dominate the engineering
schedule and make the project claim fragile to a single detection failure
\cite{deol-et-al-2024};
\cite{kienle-et-al-2025}.
Third, a real-instrument prototype would be difficult to reproduce across
sites, since each hospital has different inventory, variant mix, and
replacement cycles. Fourth, the cognitive task of interest — applying a local
tray rule to select, count, and discriminate among items — can be studied
without the sterility, weight, and handling requirements of real steel.
Printed cards preserve the selection, count, and lookalike-discrimination
demands while stripping away the confounding variables of instrument
availability, cleaning policy, and physical handling risk.
The card interface follows the same abstraction strategy as NeuroVase, which
uses tangible cue cards to teach neurovascular anatomy and stroke assessment
without requiring cadavers, imaging workstations, or clinical placements
\cite{jahani-et-al-2026-neurovase}.
QR-code-based medical education tools similarly demonstrate that low-cost
printable codes can support healthcare learning, simulation, and training
access in resource-constrained settings
\cite{karia-et-al-2019}. The AprilTag card is a deliberate methodological choice that controls for
instrument-access confounds while keeping the measured task aligned with the
real cognitive demands of tray reconstruction.
\subsection{Rationale for Gamified and Simulation-Based Learning Style}
\label{sec:6-6-rationale-for-gamified-and-simulation-based-learning-style}
TrayGuard embeds retrieval practice, feedback, scoring, progressive difficulty,
and confidence capture within a structured learning session that resembles a
serious tabletop training game. This section justifies that design against
traditional lecture-based training, video-only instruction, and unguided
self-study.
\subsubsection{Why Gamification and Simulation over Traditional Instruction}
\label{sec:6-6-1-why-gamification-and-simulation-over-traditional-instruction}
Traditional sterile-processing training relies heavily on lecture, video, and
supervised on-the-job practice \cite{aorn-staffing-shortage-2025};
\cite{hspa-crcst-accessed-2026}. Those
methods are necessary but not sufficient: lectures can introduce instrument
names and tray rules, but they do not exercise the applied discrimination and
count-sheet following that defines tray reconstruction. Simulation-based
medical education with deliberate practice has been shown to outperform
traditional clinical education in meta-analysis, and technology-enhanced
simulation improves knowledge, skills, and behaviors compared with no
intervention \cite{cook-et-al-2011};
\cite{mcgaghie-et-al-2011}. The features that
make simulation effective — feedback, repetitive practice, curriculum
integration, and measurable outcomes — map directly to TrayGuard's learning
loop \cite{issenberg-et-al-2005}.
Gamification adds a second layer. A 2021 systematic review of gamification in
health professions education found that combining assessment attributes (scores,
progress tracking) with conflict/challenge attributes (timed tasks, increasing
difficulty) increased learners' engagement with study material and led to
improved learning outcomes in most studies that included a comparison condition
\cite{van-gaalen-et-al-2021}. A more
recent umbrella review of game-based learning in nursing education found a
large pooled effect on academic knowledge performance (standardized mean
difference 1.06) and a moderate effect on skill performance (0.54)
\cite{malicki-et-al-2025}. The mechanism is
plausible: game elements increase repetition, focus attention on error patterns,
and make low-stakes practice more acceptable than high-stakes testing.
The same review
cautions that most gamified intervention studies lack well-defined control
groups and theoretical grounding for their choice of game elements
\cite{van-gaalen-et-al-2021}. TrayGuard
addresses that concern by grounding each game-like feature in a specific
learning mechanism.
\subsubsection{Specific Game Elements and Their Learning Rationale}
\label{sec:6-6-2-specific-game-elements-and-their-learning-rationale}
\begin{table}[H]
\centering
\caption{Game elements and their learning-science rationale}
\label{tab:game-elements-and-their-learning-science-rationale}
\small\renewcommand{\arraystretch}{1.2}\setlength{\tabcolsep}{4pt}
\begin{tabularx}{\textwidth}{|>{\raggedright\arraybackslash}X|>{\raggedright\arraybackslash}X|>{\raggedright\arraybackslash}X|}
\hline
Game-like feature & Learning mechanism & Evidence base \\
\hline
Pre-test establishes a baseline score & Goal-setting and self-assessment calibration before practice & Pre-testing improves later retention compared with studying without a prior test, even when answers are not provided \cite{roediger-and-karpicke-2006}; \cite{larsen-et-al-2009}. \\
\hline
Retrieval cards use prompt-before-reveal & Active recall with immediate confirmation; the core retrieval-practice effect & Tests improve long-term retention more than restudy across health-professions contexts \cite{dunlosky-et-al-2013}; \cite{barrison-et-al-2025}. \\
\hline
Quiz with score and weak-item logging & Retrieval with spaced repetition and targeted remediation & Weak-item repetition reduces error persistence on later tasks \cite{ofstead-et-al-2023}; \cite{hattie-and-timperley-2007}. \\
\hline
Practice tray sort with error-category feedback & Deliberate practice with specific, actionable correction & Immediate error-specific feedback improves skill acquisition more than correctness-only feedback \cite{hattie-and-timperley-2007}; \cite{issenberg-et-al-2005}. \\
\hline
Confidence capture after each item & Metacognitive calibration; exposes high-confidence errors & High-confidence errors are more resistant to correction and indicate overconfidence that feedback must target explicitly \cite{ofstead-et-al-2023}. \\
\hline
Parallel no-hints post-test & Near-transfer assessment without answer leakage & Separating assessment from practice prevents score inflation from memorization of exact layouts \cite{bell-et-al-2008}. \\
\hline
Metrics export and weak-item summary & Learner- and instructor-facing performance data for targeted remediation & Instructors cannot target weak items or weak rules without item-level evidence \cite{kovach-2012}; \cite{nichol-et-al-2024}. \\
\hline
\end{tabularx}
\end{table}
Each element serves a specific cognitive purpose rather than decoration.
Points and progress bars (the most common gamification elements in medical
education) are present only insofar as they support self-assessment: the
pre/post accuracy score, the weak-item recovery count, and the error-category
breakdown are the central feedback signals. The design deliberately avoids
leaderboards, badges, and competitive ranking, which can shift motivation from
intrinsic learning to extrinsic reward-seeking in health-professions contexts
\cite{van-gaalen-et-al-2021}.
\subsubsection{Why Not a More Complex Game Format}
\label{sec:6-6-3-why-not-a-more-complex-game-format}
Several alternative gamification approaches were considered and rejected.
Full virtual reality (VR) simulation offers immersive instrument handling but
requires hardware (headset, controllers), room space, onboarding time, and
technical support that the first prototype cannot assume. SteriBoost, a
commercial VR sterile-processing trainer, demonstrates the feasibility of VR
for this domain, but its hardware and custom-scene requirements place it at a
later validation stage rather than the first project prototype
\cite{steriboost-accessed-2026}.
Competitive multiplayer formats (leaderboards, team challenges) can increase
engagement but risk publicizing individual error rates in a training context,
potentially creating the punitive environment that the ethics controls
explicitly prohibit (Section~\ref{sec:12--broader-impacts-and-ethics}). They also require larger participant pools
and session coordination that the first 8-12-person study cannot support.
Simulated patient scenarios or narrative-driven learning modules (used in tools
such as Touch Surgery) can improve engagement and contextual reasoning, but
they add authoring burden and reduce the number of tray-repetition cycles per
session \cite{tulipan-et-al-2019};
\cite{medtronic-touch-surgery-accessed-2026}.
TrayGuard's target task — learning a specific local tray rule — benefits more
from repeated applied practice than from narrative context, so the game
structure stays close to the assessment loop.
The chosen gamification level is therefore lightweight, mechanism-driven, and
designed around the retrieval-practice and feedback evidence. It is closer to
the serious-game approach of SteriDefi, the SF2S sterilization serious game
that uses question banks, difficulty levels, and scoring to motivate repeated
knowledge retrieval, than to a narrative or VR simulation
\cite{vanhaverbeke-et-al-2018}.
\subsubsection{Learning Mechanisms Behind the Game Elements}
\label{sec:6-6-4-learning-mechanisms-behind-the-game-elements}
Structured, repeated practice with immediate, error-specific feedback,
delivered in a bounded session with clear goals and measurable progress,
produces measurable short-term learning. Game elements sustain the repetition
that retrieval-practice research shows is necessary for durable learning
\cite{roediger-and-karpicke-2006};
\cite{dunlosky-et-al-2013}. The "game" is the
learning loop itself — the same logic that makes flashcard apps, spaced-
repetition quizzes, and simulation-based skills trainers effective teaching
tools. A learner who improves on the TrayGuard simulated tray task has
undergone repeated retrieval practice on local names, aliases, counts, and
lookalike distinctions, followed by an applied sorting task with
category-specific feedback.
\section{Requirements Definition}
\label{sec:7--requirements-definition}
The requirements translate the narrowed problem into a formative learning
claim. They are not hospital deployment requirements. They define what the
prototype must do to support a credible project statement about novice improvement
on a simulated local tray task.
\subsection{Element Definition}
\label{sec:7-1-element-definition}
\begin{table}[H]
\centering
\caption{Element definition for the TrayGuard system}
\label{tab:element-definition-for-the-trayguard-system}
\small\renewcommand{\arraystretch}{1.2}\setlength{\tabcolsep}{4pt}
\begin{tabularx}{\textwidth}{|>{\raggedright\arraybackslash}X|>{\raggedright\arraybackslash}X|}
\hline
Element & Definition For TrayGuard \\
\hline
Intended practice & Novice preparation for local tray reconstruction before or alongside supervised SPD training. \\
\hline
Other practice affected & Instructor review, preceptor coaching, tray content updates, weak-item remediation, and future CV-supported tray review. \\
\hline
Artifact & A local training app, file-backed tray module, printable marker cards, tray-sorting surface, scoring engine, and export/report package. \\
\hline
Problem addressed & Learners need repeated, measurable practice applying local tray rules, counts, aliases, and lookalike discrimination before real tray work. \\
\hline
Technology & Local web/app workflow, JSON tray modules, QR or AprilTag cards, camera or manual selection input, scoring logic, and CSV/JSON export. \\
\hline
Primary users & Novice learners, instructors, SPD educators, and future project teams. \\
\hline
Perception goal & The system should feel like a serious practice tool, not a clinical approval device or punitive surveillance system. \\
\hline
Environment & Classroom, skills lab, project demo, or supervised training setting; not a live sterile field or production SPD deployment. \\
\hline
Core function & Convert a local tray module into study, quiz, simulated sorting, feedback, assessment, and learner metrics. \\
\hline
Required behavior & Separate practice from assessment, suppress hints during pre/post modes, classify errors consistently, and export comparable rows. \\
\hline
Structure & Content module, learning workflow, marker/manual input, scoring engine, feedback engine, reporting/export, and instructor review. \\
\hline
Intended effects & Improve simulated tray-sorting accuracy, reduce severe errors, expose weak items, calibrate confidence, and reduce early preceptor burden. \\
\hline
Possible side effects & Overclaiming simulated scores, punitive use of learner metrics, answer leakage, inequitable access to updated modules, and reduced attention to supervised hands-on practice. \\
\hline
\end{tabularx}
\end{table}
\subsection{Objective Tree}
\label{sec:7-2-objective-tree}
\begin{lstlisting}
Improve novice readiness for local tray reconstruction through repeatable,
measurable simulated practice.

1. Increase local tray knowledge
   1.1 Teach local names, aliases, and instrument families
   1.2 Teach required counts and tray-specific rules
   1.3 Teach lookalike and same-family distinctions

2. Improve applied tray-sorting performance
   2.1 Reduce missing required items
   2.2 Reduce wrong and misidentified items
   2.3 Reduce extra and wrong-count items
   2.4 Maintain reasonable task time without rewarding guessing

3. Make learning measurable
   3.1 Capture pre-test and post-test accuracy
   3.2 Export item-level error categories
   3.3 Capture confidence, Not sure choices, and high-confidence errors
   3.4 Preserve comparable assessment variants

4. Support instructors
   4.1 Show weak instruments and weak tray rules
   4.2 Keep modules local, versioned, and reviewable
   4.3 Avoid collecting patient data or unnecessary learner identifiers

5. Preserve claim boundaries
   5.1 Label outputs as simulated practice evidence
   5.2 Require expert review before content-valid claims
   5.3 Require SPD pilot evidence before workplace-effectiveness claims
\end{lstlisting}
\subsection{Performance Specifications}
\label{sec:7-3-performance-specifications}
The first build should use the following targets. They are deliberately
formative rather than clinical: they define when the prototype is ready to
support a project novice-learning claim.
\begin{table}[H]
\centering
\caption{Performance specifications and acceptance targets}
\label{tab:performance-specifications-and-acceptance-targets}
\small\renewcommand{\arraystretch}{1.2}\setlength{\tabcolsep}{4pt}
\begin{tabularx}{\textwidth}{|>{\raggedright\arraybackslash}X|>{\raggedright\arraybackslash}X|>{\raggedright\arraybackslash}X|>{\raggedright\arraybackslash}X|}
\hline
Area & Target & Evidence Basis & Measurement \\
\hline
Participant count & Minimum 8 novice participants; target 12 complete participants if schedule allows. & NeuroVase used a 40-participant controlled study for an educational card/AR system, while qualitative methods literature supports small, narrow, homogeneous samples for early formative work \cite{jahani-et-al-2026-neurovase}; \cite{malterud-et-al-2016}; \cite{hennink-and-kaiser-2022}. & Unique participants with complete pre/post attempts. \\
\hline
Completion rate & At least 7 of 8 participants, or 85\% when n > 8, complete the full loop without blocking moderator intervention. & Usability must be separated from learning; NeuroVase used SUS and user-experience measures alongside knowledge testing. & Completion flag and blocking help requests. \\
\hline
Accuracy gain & Mean paired tray accuracy improves by at least 20 percentage points from pre-test to post-test. & Ofstead's SP training pilot showed large pre/post score movement after structured teaching and hands-on practice; TrayGuard should set a smaller but visible formative threshold \cite{ofstead-et-al-2023}. & \texttt{post\_accuracy - pre\_accuracy}. \\
\hline
Post-test floor & Mean post-test accuracy is at least 80\%, and no more than one participant scores below 70\%. & The goal is not only improvement from a low baseline; learners should reach a usable simulated familiarity level before the module is considered ready. & Post-test \texttt{accuracy\_score}. \\
\hline
Severe-error reduction & Mean \texttt{missing + wrong + misidentified} errors decrease by at least 30\%. & Alfred et al. and Nichol et al. show missing, wrong, extra, and assembly-related errors are operationally meaningful categories \cite{alfred-et-al-2021}; \cite{nichol-et-al-2024}. & Paired error-category counts. \\
\hline
Speed without unsafe guessing & Median post-test duration is at least 10\% lower, or is no more than 20\% higher if accuracy improves by at least 20 points. & Time matters because tray assembly is workflow-bound, but speed alone could reward guessing. & \texttt{duration\_seconds}, paired by learner. \\
\hline
Confidence calibration & High-confidence errors decline by at least 30\%; low-confidence correct answers are exported. & Ofstead captured confidence/satisfaction, and feedback research supports showing learners where they are and what to do next \cite{ofstead-et-al-2023}; \cite{hattie-and-timperley-2007}. & Confidence >= 4 on incorrect items; confidence <= 2 on correct items. \\
\hline
Retention & A 7-day delayed tray sort retains at least half of the immediate accuracy gain and delayed accuracy does not return to within 5 percentage points of the pre-test baseline. & Bell et al. showed immediate medical-learning gains can decay within days and recommended reinforcement after as little as 1 week; same-day results should not be called retention \cite{bell-et-al-2008}. The 5pp falsification floor ensures retention is measured against baseline rather than as any positive number. & Delayed \texttt{accuracy\_score} relative to pre/post. \\
\hline
Convergent evidence & Post-test tray-sort accuracy correlates at r >= 0.50 with an independent photo-identification test of the same instruments administered in the same session. & Messick's external-aspect and Cook et al.'s relations-to-other-variables standard require that scores on the target task co-vary with a separate measure of the same construct \cite{messick-validity-1989, cook-beckman-validity-2006}. A photo-ID test shares the knowledge domain (instrument identity, aliases, lookalike discrimination) while differing in response format (identification without tray-sort context), so a moderate-to-strong correlation supports convergent validity. NeuroVase provides a recent precedent: Jahani et al. used independent knowledge tests alongside tangible-card simulation performance and found moderate correlations between measures \cite{jahani-et-al-2026-neurovase}. & Pearson r between post-test tray accuracy and photo-ID accuracy across participants. \\
\hline
Marker reliability & At least 95\% card detection in normal layouts and no systematic marker-ID confusion before user testing. & The marker layer is a measurement tool; if detection fails, scores confound learner error with app error. & Marker bench-test log. \\
\hline
Export completeness & Every completed run exports attempt-level and item-level records with required fields. & Instructor-facing evidence is part of the training value proposition. & CSV/JSON schema validation. \\
\hline
Demo success & A facilitator can load the module, run pre-test, study, quiz, practice sort, post-test, and export metrics in one uninterrupted demo. & The project requires an interactive demonstration of a functional design. & End-to-end checklist. \\
\hline
\end{tabularx}
\end{table}
\subsection{Quality Function Deployment}
\label{sec:7-4-quality-function-deployment}
Compact QFD:
\begin{table}[H]
\centering
\caption{Quality function deployment: needs to engineering targets}
\label{tab:quality-function-deployment-needs-to-engineering-targets}
\small\renewcommand{\arraystretch}{1.2}\setlength{\tabcolsep}{4pt}
\begin{tabularx}{\textwidth}{|>{\raggedright\arraybackslash}X|>{\raggedright\arraybackslash}X|>{\raggedright\arraybackslash}X|>{\raggedright\arraybackslash}X|}
\hline
Stakeholder need & Engineering characteristic & Target value & Verification \\
\hline
Learners can practice local tray content repeatedly. & File-backed module with required items, distractors, aliases, features, counts, and variants. & One complete \texttt{basic\_general\_tray\_v1} module with 8-12 required concepts and 3-6 distractors. & Module schema review and demo load. \\
\hline
Practice feels connected to tray reconstruction, not isolated memorization. & Simulated tray sort using physical cards or manual selection. & Pre-test, practice, post-test, and optional delayed variant all use the same tray rule. & End-to-end run. \\
\hline
Educators can see what learners missed. & Attempt-level and item-level export. & 100\% of completed runs export all required fields. & Schema validation. \\
\hline
Errors map to operationally meaningful categories. & Scoring engine with \texttt{missing}, \texttt{extra}, \texttt{wrong}, \texttt{misidentified}, and \texttt{wrong\_count}. & Every submitted tray receives category counts. & Unit tests or scored example cases. \\
\hline
Assessment is not contaminated by coaching. & Mode boundaries and feedback suppression. & No hints or corrective feedback in pre-test, post-test, or delayed modes. & UI walkthrough and export mode check. \\
\hline
The physical marker layer does not corrupt scores. & Marker reliability and user confirmation. & At least 95\% detection in normal layouts; confirmation before scoring. & Study 0 marker log. \\
\hline
Learner confidence is visible. & Confidence or \texttt{Not sure} capture. & High-confidence errors and low-confidence correct answers exported. & Export inspection. \\
\hline
The system remains ethical and nonpunitive. & Privacy and report labels. & Hashed learner IDs, no patient data, and report language that says practice evidence. & Report review. \\
\hline
Future teams can extend the design. & Versioned data model and traceability. & Module version, source notes, verification status, and assessment variant IDs included. & Data file review. \\
\hline
\end{tabularx}
\end{table}
\subsection{Requirement Traceability Matrix}
\label{sec:7-5-requirement-traceability-matrix}
\begin{table}[H]
\centering
\caption{Requirement traceability matrix}
\label{tab:requirement-traceability-matrix}
\small\renewcommand{\arraystretch}{1.15}\setlength{\tabcolsep}{3pt}
\begin{tabularx}{\textwidth}{|>{\raggedright\arraybackslash}p{1.5cm}|>{\raggedright\arraybackslash}X|>{\raggedright\arraybackslash}X|>{\raggedright\arraybackslash}p{2.5cm}|>{\raggedright\arraybackslash}p{2.3cm}|>{\raggedright\arraybackslash}p{1.7cm}|}
\hline
Req. & Stakeholder need & Design feature & Evidence source & Validation method & Status \\
\hline
FR-1 Load local module & Educators need updateable local tray rules. & JSON tray module and module selector. & Count-sheet and local tray evidence \cite{nadeau-2024}. & Change module data and reload. & Ready for impl. \\
\hline
FR-2 Retrieval cards & Learners need active recall of names, features, and counts. & Prompt-before-reveal cards. & Retrieval and flashcard evidence \cite{roediger-and-karpicke-2006}; \cite{barrison-et-al-2025}. & Card walkthrough and quiz logs. & Ready for impl. \\
\hline
FR-3 Pre/post assessment & Reviewers need measurable learning evidence. & No-hints pre-test and post-test variants. & Ofstead pre/post pattern \cite{ofstead-et-al-2023}. & Paired run export. & Protocol-ready. \\
\hline
FR-4 Quiz mode & Learners need feedback before applied sorting. & Quiz with confidence and weak-item logging. & Feedback and retrieval evidence \cite{hattie-and-timperley-2007}. & Quiz result export. & Ready for impl. \\
\hline
FR-5 Practice feedback & Learners need actionable correction. & Error-specific feedback engine. & Tray defect categories \cite{alfred-et-al-2021}; \cite{nichol-et-al-2024}. & Seeded error cases. & Defined. \\
\hline
FR-6 Metrics export & Instructors need reviewable evidence. & Attempt and item CSV/JSON. & Error-reporting and competency evidence \cite{nichol-et-al-2024}; \cite{ofstead-et-al-2023}. & Schema validation. & Defined. \\
\hline
FR-7 Mode separation & Assessment must remain fair. & Mode state and feedback suppression. & Assessment design and answer-leakage risk. & UI walkthrough. & Defined. \\
\hline
NFR-1 Fast data load & Demo cannot stall. & Local file storage. & Usability expectations. & Load timing <= 2 seconds. & Future test. \\
\hline
NFR-2 Short flow & Participants can finish in one session. & Bounded module and timed segments. & Study protocol and sample feasibility. & Timing <= 20 min for demo. & Future test. \\
\hline
NFR-3 Export reliability & No completed run should be lost. & Local export after each run. & Study evidence requirement. & 100\% completed-run exports. & Future test. \\
\hline
NFR-4 Privacy & Learners should not be surveilled unnecessarily. & Hashed IDs and no PHI. & Healthcare data governance concerns. & Export review. & Defined. \\
\hline
NFR-5 Content traceability & Local content must be auditable. & Source notes and instructor verification status. & Count-sheet and local content boundary. & Module review. & Defined. \\
\hline
NFR-6 Capability labels & Stakeholders must not overread scores. & Report language: simulated practice evidence. & Certification and clinical-transfer boundaries. & Report review. & Defined. \\
\hline
\end{tabularx}
\end{table}
\section{System Design}
\label{sec:8--system-design}
The product's system design follows a documented arc: the original
camera-supported tray-check plan, the CV experiments that showed why that plan
did not generalize, and the current training-platform architecture.

\subsection{Original Planned System: Camera-Supported Tray Check}
\label{sec:8-1-original-planned-system-camera-supported-tray-check}

The original product concept was a tray-verification tool that used a camera,
a YOLO object detector, and confidence-based logic to check whether a
technician's assembled tray matched the required instrument list. The system
was designed to support the following workflow:

\begin{enumerate}
  \item Load or enter the required instrument list for a tray.
  \item Capture the current tray view with a camera.
  \item Detect visible instruments and assign class labels with confidence scores.
  \item Aggregate detections into observed class counts.
  \item Compare observed counts against required counts.
  \item Present a result that highlights what appears present, missing, extra, or uncertain.
  \item Preserve the final decision as human judgment rather than silent automation.
\end{enumerate}

The tray-status logic used a two-threshold confidence band with a scan-quality
gate, because a single detector threshold is insufficient to distinguish a
truly absent item from one missed due to glare, occlusion, blur, similar-class
confusion, or open-set behavior:

\begin{table}[H]
\centering
\caption{Original planned tray-status logic using two-threshold confidence band}
\label{tab:original-planned-tray-status-logic}
\small\renewcommand{\arraystretch}{1.2}\setlength{\tabcolsep}{4pt}
\begin{tabularx}{\textwidth}{|>{\raggedright\arraybackslash}X|>{\raggedright\arraybackslash}X|}
\hline
User-Facing State & Decision Logic \\
\hline
\texttt{Present} & Required class has at least required count of detections above high threshold \texttt{T\_present}, no active similar-class or image-quality risk flag, and duplicate boxes reconciled. \\
\hline
\texttt{Needs Review} & Detections between \texttt{T\_review} and \texttt{T\_present}, similar class appears, count unstable, unknown or extra object plausible, or scan quality degraded. \\
\hline
\texttt{Missing} & No candidate above low threshold \texttt{T\_review}, and scan quality acceptable enough that absence is meaningful. \\
\hline
\texttt{Rescan Recommended} & System lacks evidence but glare, occlusion, blur, poor framing, or lighting makes a missing label unsafe. \\
\hline
\end{tabularx}
\end{table}

The technician-facing UI was designed around the tray decision with one
primary action per state, visual-first feedback, alert restraint, and explicit
human sign-off:

\begin{table}[H]
\centering
\caption{Original planned technician UI requirements}
\label{tab:original-planned-technician-ui-requirements}
\small\renewcommand{\arraystretch}{1.2}\setlength{\tabcolsep}{4pt}
\begin{tabularx}{\textwidth}{|>{\raggedright\arraybackslash}X|>{\raggedright\arraybackslash}X|}
\hline
Requirement & Design Direction \\
\hline
One primary action per state & Show the next meaningful action: start scan, rescan, review issue, confirm tray, or stop. \\
\hline
Clear tray-status hierarchy & Summarize Complete, Missing, Extra, and Needs Review above detailed detections. \\
\hline
Fast correction and rescan & Provide direct correction, dismiss, and rescan paths instead of forcing a workflow restart. \\
\hline
Specific review reasons & Label review prompts with the cause: low confidence, similar-class risk, glare, overlap, unknown object, or count mismatch. \\
\hline
Visual-first feedback with optional sound & Persistent visible status as the source of truth; short optional audio for major state changes. \\
\hline
Alert restraint & Escalate only conditions that change the tray decision or require user action. \\
\hline
Human sign-off & Require explicit final confirmation. The system supports the decision, not silently makes it. \\
\hline
Manual fallback & Make it clear how to continue if camera, model, or scan quality is unavailable. \\
\hline
\end{tabularx}
\end{table}

The original system also included a physical capture fixture: an overhead
camera stand mounted on a portable weighted board, designed for tabletop use
with an optional rolling standing base. The board-mounted design preserved
repeatable camera geometry across work surfaces without requiring table
modification.

\subsection{Why the Original Design Was Revised}
\label{sec:8-2-why-the-original-design-was-revised}

The CV experiments (Section~\ref{sec:6-2-1-what-the-cv-experiments-showed})
showed that detection degraded under every real-world variation tested:
lighting changes caused mAP50-90 decay from separated to overlay conditions,
shape-similarity experiments produced misidentified predictions above 30\%
within same-family pairs, and open-set tests showed false-positive rates of
40--60\% for unknown objects. The original camera-based tray checker could
not generalize across the environments, instrument sets, and lighting
conditions that a real SPD deployment would require.

The project therefore pivoted from a CV-based verification tool to a training
platform that uses simulated tray sorting. The camera infrastructure remains
available for future support-layer work, but the product's primary evidence
path no longer depends on real-instrument detection.

\subsection{Revised System: Training Platform}
\label{sec:8-3-revised-system-training-platform}

The selected system is a local, file-backed training app with a physical card
interface. It should be implemented as a buildable learning loop first, with
real-instrument CV reserved for later support work.
\subsection{System Overview}
\label{sec:8-3-1-system-overview}
TrayGuard should be built as a local module runner plus a card-scanning
assessment surface. The app does not need to recognize real instruments for the
first study. It needs to observe which simulated instrument cards the learner
placed in the tray area, score that selection against a local tray template,
and export evidence.
\begin{lstlisting}
local tray module
  -> pre-test AprilTag tray sort
  -> retrieval-first study cards
  -> quiz with confidence / Not sure
  -> practice AprilTag tray sort with feedback
  -> post-test AprilTag tray sort
  -> metrics export
   -> committed 7-day delayed retention sort (see falsification criteria in Section~\ref{sec:10-8-study-2-delayed-retention-check})
\end{lstlisting}
The architecture follows the evidence base:
\begin{itemize}
  \item Local tray modules are required because count sheets define tray names, contents, quantities, sizes, and reference/catalog numbers, and real count sheets include fields such as reference number, description, location, quantity, check box, total count, and hospital signature \cite{nadeau-2024}; \cite{stryker-gamma4-count-sheet-2023}; \cite{stryker-imn-count-sheet-2023}.
  \item Retrieval-first study and quiz modes are justified by retrieval practice, health-professions flashcard evidence, and spaced digital education \cite{roediger-and-karpicke-2006}; \cite{barrison-et-al-2025}; \cite{martinengo-et-al-2024}.
  \item Pre/post practice, confidence capture, and optional delayed retention follow the sterile-processing precedent in Ofstead and the retention warning in Bell et al. \cite{ofstead-et-al-2023}; \cite{bell-et-al-2008}.
   \item AprilTag cards are a physical proxy, not a real-instrument detector. The tangible-card idea is supported by NeuroVase-style cue cards and board/card-based healthcare education, but clinical transfer remains future validation \cite{jahani-et-al-2026-neurovase}; \cite{ward-et-al-2019}.
\end{itemize}
\subsection{Sensing Architecture Alternatives}
\label{sec:8-3-2-sensing-architecture-alternatives}
The system must observe which instruments the learner selects. Three sensing
architectures were evaluated before selecting the marker-card approach:
\begin{table}[H]
\centering
\caption{Sensing architecture candidates and selection rationale}
\label{tab:sensing-architecture-candidates-and-selection-rationale}
\small\renewcommand{\arraystretch}{1.2}\setlength{\tabcolsep}{4pt}
\begin{tabularx}{\textwidth}{|>{\raggedright\arraybackslash}X|>{\raggedright\arraybackslash}X|>{\raggedright\arraybackslash}X|>{\raggedright\arraybackslash}X|}
\hline
Architecture & Description & Advantages & Main Limits \\
\hline
Workspace observation & Camera and CV model detect visible instruments, count known classes, compare against required list. & Commodity hardware; matches the visual-verification failure mode. & Sensitive to lighting, reflectivity, occlusion, similar-class confusion, and open-set objects. \\
\hline
Tagged identification & RFID tags, barcodes, or engraved identifiers read during tray preparation. & Strong item-level identity; less dependent on visual conditions. & Requires tags, readers, scanner workflows, and maintenance; infrastructure and workflow costs are documented barriers \cite{olivere-et-al-2021}; \cite{kusuda-et-al-2024}; \cite{coustasse-et-al-2013}. \\
\hline
Instrumented placement & Custom tray slots or embedded sensors confirm occupied locations. & Simple placement checks in a standardized tray. & Requires specialized trays, restricts layout flexibility, and cannot discriminate wrong-but-similar instruments. \\
\hline
\end{tabularx}
\end{table}
Workspace observation is the strongest long-term option because it matches the
existing camera infrastructure and the visual nature of the task. Tagged
identification and instrumented placement impose hardware and workflow costs
that raised the project's adoption burden. The selected marker-card approach is
a simulation-layer compromise: it preserves physical search, selection, and
count-sheet following without requiring either deployment-grade CV or
instrument-level hardware. The CV code remains available as future support for
workspace observation once deployment validation is feasible.
\subsection{Subsystem Breakdown}
\label{sec:8-3-3-subsystem-breakdown}
The system separates into these functional subsystems:
\begin{table}[H]
\centering
\caption{Functional subsystem breakdown}
\label{tab:functional-subsystem-breakdown}
\small\renewcommand{\arraystretch}{1.2}\setlength{\tabcolsep}{4pt}
\begin{tabularx}{\textwidth}{|>{\raggedright\arraybackslash}X|>{\raggedright\arraybackslash}X|}
\hline
Subsystem & Responsibility \\
\hline
Local content/tray module & Stores instrument records, aliases, photos, distinguishing features, required quantities, distractors, and tray-template versions. \\
\hline
Learning workflow & Routes the learner through pre-test, study, quiz, practice sort, weak-item review, post-test, and export. \\
\hline
Retrieval practice & Presents prompt-before-answer cards and quizzes for instrument names, features, lookalikes, and required counts. \\
\hline
AprilTag card scanning & Detects printed marker cards and maps marker IDs to instrument IDs so the app can observe a physical tray simulation. \\
\hline
Simulated tray scoring & Compares detected cards and selected quantities against the tray template. \\
\hline
Feedback engine & Reports missing, extra, wrong, misidentified, and wrong-count errors in practice mode. \\
\hline
Confidence/uncertainty capture & Records confidence or \texttt{Not sure} responses for assessment and instructor review. \\
\hline
Reporting/export & Produces pre/post accuracy, time, error-category, confidence, and weak-item summaries. \\
\hline
Instructor review & Lets an educator verify tray content, local naming, distractors, and study evidence before broader use. \\
\hline
\end{tabularx}
\end{table}
Computer vision for recognizing real instruments remains a future support
subsystem. In the project build, the marker-scanning subsystem is deliberately a
proxy for observing user decisions, not a claim that the app can identify real
surgical instruments under clinical conditions.
\subsection{Functional Block Diagram}
\label{sec:8-3-4-functional-block-diagram}
\begin{lstlisting}
[Mermaid diagram]
flowchart LR
  A["Local Tray Module JSON"] --> B["Learning Workflow Controller"]
  B --> C["Study Cards"]
  B --> D["Quiz"]
  B --> E["AprilTag Tray Sort"]
  E --> F["Detected Cards"]
  F --> G["User Confirmation"]
  G --> H["Scoring Engine"]
  A --> H
  H --> I["Practice Feedback"]
  H --> J["Assessment Result"]
  J --> K["Metrics Export"]
  K --> L["Learner / Instructor Summary"]
\end{lstlisting}
In practice modes, the scoring engine sends error-specific feedback back to the
learner. In assessment modes, the same engine records errors but suppresses
corrective feedback until after submission.
\subsection{Dependency Diagram}
\label{sec:8-3-5-dependency-diagram}
\begin{lstlisting}
[Mermaid diagram]
flowchart TD
  A["Instructor-verified local tray module"] --> B["Retrieval cards and quiz"]
  A --> C["Assessment variants"]
  A --> D["Scoring engine"]
  E["Printable neutral marker cards"] --> F["AprilTag / QR detection"]
  F --> G["User confirmation"]
  G --> D
  C --> H["Pre/post/delayed tray sort"]
  B --> I["Practice and weak-item review"]
  I --> H
  D --> J["Attempt and item export"]
  J --> K["Learner / instructor report"]
  L["Study protocol"] --> H
  L --> K
  M["Expert review"] --> A
  M --> K
\end{lstlisting}
Highest-risk dependency wires:
\begin{table}[H]
\centering
\caption{Highest-risk dependency wires and mitigations}
\label{tab:highest-risk-dependency-wires-and-mitigations}
\small\renewcommand{\arraystretch}{1.2}\setlength{\tabcolsep}{4pt}
\begin{tabularx}{\textwidth}{|>{\raggedright\arraybackslash}X|>{\raggedright\arraybackslash}X|>{\raggedright\arraybackslash}X|}
\hline
Dependency & Risk & Mitigation \\
\hline
Expert review -> local tray module & If names, counts, distractors, or feedback are unrealistic, learning results lose content validity. & Require instructor/SPD review before claiming content plausibility. \\
\hline
Marker cards -> detection -> scoring & If cards are missed or confused, learner scores mix app error with user error. & Run Study 0, log raw detections, and require user confirmation before scoring. \\
\hline
Assessment variants -> pre/post claim & If the post-test is easier or leaks answers, improvement is not credible. & Match required items, counts, distractors, lookalikes, and difficulty across variants. \\
\hline
Retrieval/practice -> post-test & If practice only teaches the exact post-test layout, results are memorization. & Use parallel variants and test the same tray rule in a different order/view/distractor mix. \\
\hline
Export -> instructor report & If exported rows are incomplete, results cannot support the project claim. & Validate schemas before running participants. \\
\hline
Student sample -> stakeholder conclusion & Students can test novice learning, but not SPD workplace transfer. & Label Study 1 as simulated novice evidence and reserve SPD claims for later pilot. \\
\hline
\end{tabularx}
\end{table}
\subsection{Data Model And Module Contract}
\label{sec:8-3-6-data-model-and-module-contract}
The data model should mirror a real count sheet while staying small enough for
the first build. Stryker's public count sheets show the practical minimum:
reference number, description, location/layer, quantity, check box, total count,
additional items, and hospital signature. TrayGuard adds learning-specific
fields for aliases, lookalikes, confidence, and error categories.
\begin{table}[H]
\centering
\caption{Data model objects and required fields}
\label{tab:data-model-objects-and-required-fields}
\small\renewcommand{\arraystretch}{1.2}\setlength{\tabcolsep}{4pt}
\begin{tabularx}{\textwidth}{|>{\raggedright\arraybackslash}X|>{\raggedright\arraybackslash}X|>{\raggedright\arraybackslash}X|}
\hline
Object & Required Fields & Why It Matters \\
\hline
\texttt{Instrument} & \texttt{id}, \texttt{display\_name}, \texttt{family}, \texttt{aliases}, \texttt{distinguishing\_features}, \texttt{image\_refs}, \texttt{apriltag\_id}, \texttt{local\_verification\_status} & Supports retrieval cards, quiz prompts, card scanning, and local naming. \\
\hline
\texttt{TrayTemplate} & \texttt{id}, \texttt{name}, \texttt{version}, \texttt{procedure\_family}, \texttt{source\_note}, \texttt{required\_items}, \texttt{distractor\_items}, \texttt{lookalike\_pairs}, \texttt{assessment\_variants} & Stores the local count-sheet task and separates required instruments from distractors. \\
\hline
\texttt{TrayItemRequirement} & \texttt{instrument\_id}, \texttt{quantity}, \texttt{reference\_number}, \texttt{location\_or\_layer}, \texttt{acceptable\_substitutes}, \texttt{notes} & Mirrors real count-sheet fields and makes count/location errors measurable. \\
\hline
\texttt{AssessmentVariant} & \texttt{id}, \texttt{mode}, \texttt{random\_seed}, \texttt{required\_item\_ids}, \texttt{distractor\_item\_ids}, \texttt{card\_front\_policy}, \texttt{feedback\_enabled} & Keeps pre/post/delayed tasks comparable while preventing answer leakage. \\
\hline
\texttt{LearningRun} & \texttt{run\_id}, \texttt{learner\_id\_hash}, \texttt{participant\_group}, \texttt{tray\_template\_id}, \texttt{variant\_id}, \texttt{mode}, \texttt{started\_at}, \texttt{completed\_at} & Separates study, quiz, practice, pre-test, post-test, and delayed attempts. \\
\hline
\texttt{DetectedCard} & \texttt{run\_id}, \texttt{apriltag\_id}, \texttt{instrument\_id}, \texttt{timestamp}, \texttt{detection\_confidence}, \texttt{confirmed\_by\_user} & Lets marker failures be audited instead of silently counted as learner errors. \\
\hline
\texttt{AttemptResult} & \texttt{run\_id}, \texttt{selected\_items}, \texttt{duration\_seconds}, \texttt{accuracy\_score}, \texttt{required\_recall}, \texttt{errors}, \texttt{confidence\_summary} & Provides the main pre/post evidence. \\
\hline
\texttt{ItemResult} & \texttt{run\_id}, \texttt{instrument\_id}, \texttt{required\_qty}, \texttt{selected\_qty}, \texttt{error\_category}, \texttt{confidence}, \texttt{not\_sure} & Supports item-level feedback, weak-item review, high-confidence error reporting, and instructor review. \\
\hline
\end{tabularx}
\end{table}
The first implementation can store these objects in JSON files and export CSV
for analysis. The important constraint is schema stability: pre-test, post-test,
and delayed retention must produce comparable rows.
Sample instrument JSON:
\begin{lstlisting}
{
  "id": "mayo_scissors_straight_55",
  "display_name": "Straight Mayo scissors, 5.5 in",
  "family": "cutting_dissecting",
  "aliases": [
    "straight Mayo",
    "suture scissors"
  ],
  "distinguishing_features": [
    "Heavy straight blades",
    "Broader blade profile than Metzenbaum scissors",
    "Ring-handled cutting instrument"
  ],
  "common_confusions": [
    "mayo_scissors_curved_55",
    "metzenbaum_scissors_curved_55"
  ],
  "image_refs": [
    {
      "id": "view_a",
      "path": "data/instruments/mayo_scissors_straight_55_view_a.jpg",
      "source": "local_photo",
      "approved_for_study": true,
      "approved_for_assessment": true,
      "attribution": "Local TrayGuard photo"
    },
    {
      "id": "card_front_neutral",
      "path": "outputs/cards/mayo_scissors_straight_55_neutral_front.png",
      "source": "generated_card",
      "approved_for_study": false,
      "approved_for_assessment": true,
      "attribution": "TrayGuard AprilTag card"
    }
  ],
  "apriltag_id": 12,
  "card_front_policy": "neutral_assessment_front",
  "study_prompts": [
    {
      "id": "name_recall",
      "prompt": "Name this instrument.",
      "answer": "Straight Mayo scissors, 5.5 in"
    },
    {
      "id": "feature_recall",
      "prompt": "What feature helps distinguish this from Metzenbaum scissors?",
      "answer": "Straight Mayo scissors have heavier, broader blades."
    }
  ],
  "notes": "Use local photos before pilot. Do not rely on vendor images for assessed study content.",
  "local_verification_status": "pending_review"
}
\end{lstlisting}
\subsection{Mode Behavior}
\label{sec:8-3-7-mode-behavior}
\begin{table}[H]
\centering
\caption{Mode behavior specification across all learning modes}
\label{tab:mode-behavior-specification-across-all-learning-modes}
\small\renewcommand{\arraystretch}{1.2}\setlength{\tabcolsep}{4pt}
\begin{tabularx}{\textwidth}{|>{\raggedright\arraybackslash}X|>{\raggedright\arraybackslash}X|>{\raggedright\arraybackslash}X|>{\raggedright\arraybackslash}X|>{\raggedright\arraybackslash}X|}
\hline
Mode & Input & Hints & Feedback & Stored Output \\
\hline
Pre-test & AprilTag tray sort against count-sheet prompt & No & No & Attempt result, item results, confidence, duration, errors. \\
\hline
Study cards & Instrument prompt, photo/card, feature, alias, or count & Yes after reveal & Yes & Optional viewed cards and self-rated confidence. \\
\hline
Quiz & Recall or recognition prompt with distractors & Limited after answer & Yes & Correctness, duration, confidence, weak items. \\
\hline
Practice sort & AprilTag tray sort with distractors & Yes & Yes & Attempt result plus immediate missing/extra/wrong/misidentified/wrong-count feedback. \\
\hline
Post-test & Parallel AprilTag tray sort & No & No & Comparable assessment result. \\
\hline
Delayed retention & Third parallel AprilTag tray sort after 7 days & No & No & Retention result and decay/recovery summary (with falsification criteria). \\
\hline
Export/report & Completed run data & N/A & N/A & CSV/JSON plus learner summary. \\
\hline
\end{tabularx}
\end{table}
\subsection{User Interaction / Experience}
\label{sec:8-3-8-user-interaction-experience}
Learner workflow:
\begin{enumerate}
  \item Enter anonymous participant ID.
  \item Complete prior-experience questions.
  \item Assemble the pre-test tray from neutral AprilTag instrument cards.
  \item Confirm detected cards before submission.
  \item Study retrieval-first cards and answer quiz prompts.
  \item Practice tray sorting with feedback.
  \item Complete a no-hints post-test and short survey.
  \item Return for delayed retention if scheduled.
\end{enumerate}
Instructor/researcher workflow:
\begin{enumerate}
  \item Author or load the tray module.
  \item Print neutral assessment cards and study/reference cards.
  \item Run marker reliability check.
  \item Moderate sessions and record help requests.
  \item Export pre/post/delayed metrics.
  \item Review weak items, high-confidence errors, scan failures, and qualitative feedback.
  \item Revise module content before future SPD validation.
\end{enumerate}
\section{Training Module Design}
\label{sec:9--training-module-design}
The first module is one evidence-complete local tray curriculum, not a broad
instrument encyclopedia. The module must be small enough to build, test, and
explain in a project demo while still exercising the target task: local tray
assembly from a count-sheet-like rule.
The bounded capability claim is:
\begin{quote}
TrayGuard gives novices repeated, measured practice on local tray familiarity
and can test whether that practice improves simulated tray-sorting
performance.
\end{quote}
\subsection{Module Scope}
\label{sec:9-1-module-scope}
The first module should contain:
\begin{itemize}
  \item 8-12 required instrument concepts;
  \item 3-6 distractor instruments;
  \item 2-4 known lookalike or same-family pairs;
  \item local names, aliases, family labels, required counts, and distinguishing features;
  \item AprilTag cards that represent instruments during the physical tray simulation;
  \item parallel pre-test, post-test, and delayed-retention tray variants with matched difficulty;
  \item instructor or expert review before results are presented as content-valid.
\end{itemize}
The cards should use neutral assessment fronts. Names, counts, and obvious
answer labels should be hidden during pre-test, post-test, and retention tasks
so the participant cannot solve the task by reading the card.
\subsection{FGVC-Derived Module Plausibility}
\label{sec:9-1-1-fgvc-derived-module-plausibility}
The FGVC-derived modules are practice modules built from the major-tray image
classes reported by Atabuzzaman et al. \cite{atabuzzaman-et-al-2025}. The
source contributes instrument names and images for a bounded image-card set:
scissors, clamps, forceps, needle holders, retractors, suction, and other
classes present in the major-tray archive. It gives TrayGuard a coherent pool of
real instrument classes photographed under one dataset program.

Plausibility for these modules comes from three linked sources of evidence.
First, every required card, distractor card, and lookalike pair is selected from
the FGVC major-tray class mapping. Second, the module schema mirrors real
tray/count-sheet fields: tray name, required item, quantity, reference number,
location or layer, total count, and check status
\cite{stryker-gamma4-count-sheet-2023};
\cite{stryker-imn-count-sheet-2023}. Third, the module size follows the project
study constraint: 8-12 required concepts, 3-6 distractors, and 2-4 lookalike
pairs, so novices can complete pre-test, practice, post-test, and delayed
retention tasks in a single session. The FGVC-derived quantities are simulated
as one card per selected class. Expert review sets the final clinical wording,
counts, substitutes, and local tray membership before an SPD pilot uses the
module as content-valid tray material.
\subsection{Real Tray List Examples}
\label{sec:9-1-2-real-tray-list-examples}
The first TrayGuard module should be smaller than a clinical tray, but its data
fields should look like real tray/count-sheet data. Three examples are useful:
\begin{table}[H]
\centering
\caption{Real tray list examples and design implications}
\label{tab:real-tray-list-examples-and-design-implications}
\small\renewcommand{\arraystretch}{1.2}\setlength{\tabcolsep}{4pt}
\begin{tabularx}{\textwidth}{|>{\raggedright\arraybackslash}X|>{\raggedright\arraybackslash}X|>{\raggedright\arraybackslash}X|}
\hline
Real Example & What It Shows & Design Implication \\
\hline
Stryker Gamma4 Indication Tray count sheet & A real vendor count sheet uses tray name, insert/base sections, reference numbers, descriptions, locations, quantities, check boxes, total instrument counts, additional items, and hospital signature \cite{stryker-gamma4-count-sheet-2023}. & TrayGuard should store \texttt{reference\_number}, \texttt{description/display\_name}, \texttt{location\_or\_layer}, \texttt{quantity}, and \texttt{additional\_items} fields even if the first module only uses names and quantities. \\
\hline
Stryker IMN Basic Instruments Tray count sheet & A second real count sheet uses the same pattern and splits the tray into insert/base groupings with total counts of 15 and 10 instruments \cite{stryker-imn-count-sheet-2023}. & The data model should support sections/layers and total-count validation, not just a flat list. \\
\hline
Published major orthopedic tray study & Toor et al. observed an 88-instrument major orthopedic tray across 80 procedures and compared optimized tray configurations of 47, 67, and 51 instruments \cite{toor-et-al-2022}. & Real clinical trays can be far larger than the first module. The first study should use 8-12 required concepts for feasibility, while preserving fields that can scale to larger trays. \\
\hline
\end{tabularx}
\end{table}
For the first build, use a teaching-scale tray: 8-12 required concepts, 3-6
distractors, and 2-4 lookalike pairs. This is intentionally smaller than the
real examples so a novice can complete the loop in 45-60 minutes.
\subsection{Learning Loop}
\label{sec:9-2-learning-loop}
The product is one learning loop:
\begin{lstlisting}
local tray module -> pre-test -> retrieval-first cards -> quiz -> practice tray sort -> post-test -> metrics export
\end{lstlisting}
Retrieval practice is the instructional engine. The learner should usually
commit to an answer before the system reveals it: name the instrument, identify
a distinguishing feature, choose the lookalike, recall the count, or decide
whether the item belongs in the tray. The simulated tray sort is applied
retrieval: the learner must use those same facts in context with distractors
and quantities.

Each feature in the learning loop maps to a specific evidence-based role:
\begin{table}[H]
\centering
\caption{Feature-to-evidence mapping for the learning loop}
\label{tab:feature-to-evidence-mapping-for-the-learning-loop}
\small\renewcommand{\arraystretch}{1.2}\setlength{\tabcolsep}{4pt}
\begin{tabularx}{\textwidth}{|>{\raggedright\arraybackslash}X|>{\raggedright\arraybackslash}X|>{\raggedright\arraybackslash}X|}
\hline
Feature & Evidence-based role & Engineering requirement \\
\hline
Retrieval-first study cards & Initial instruction that still asks the learner to answer before reveal & Store photos, names, aliases, functions, distinguishing features, and local notes; present prompts before explanations \\
\hline
Identification quiz & Central retrieval practice for instrument identity and features & Ask recall or recognition questions, randomize distractors, record correctness, confidence, and time \\
\hline
Simulated tray sorting & Safe applied retrieval in a realistic work task & Let learners assemble or check a tray against a count sheet in a low-risk training environment \\
\hline
Immediate error-specific feedback & Deliberate practice and error correction & Explain wrong answers and classify errors as missing, extra, wrong, misidentified, or wrong count \\
\hline
Confidence ratings & Metacognitive signal for learner and instructor & Capture confidence before or after tasks so low-confidence correct answers and high-confidence errors are visible \\
\hline
Separated pre/post tests & Evaluation not contaminated by coaching & Disable hints and feedback during tests; export comparable pre/post accuracy, time, error type, and confidence \\
\hline
\end{tabularx}
\end{table}
\subsection{Data Contract}
\label{sec:9-3-data-contract}
Use the module contract in Section~\ref{sec:8-3-6-data-model-and-module-contract} as the implementation source of truth.
The key requirement is that \texttt{Instrument}, \texttt{TrayTemplate},
\texttt{TrayItemRequirement}, \texttt{AssessmentVariant}, \texttt{LearningRun}, \texttt{DetectedCard},
\texttt{AttemptResult}, and \texttt{ItemResult} are stable enough to export comparable
pre-test, post-test, and delayed-retention rows.
Scoring should classify errors as:
\begin{itemize}
  \item \texttt{missing}: required item not selected;
  \item \texttt{extra}: non-required item selected;
  \item \texttt{wrong}: required slot filled with the wrong instrument;
  \item \texttt{misidentified}: selected item is a known lookalike or distractor for the intended item;
  \item \texttt{wrong\_count}: correct item selected with the wrong quantity.
\end{itemize}
\subsection{Mode Behavior}
\label{sec:9-4-mode-behavior}
\begin{table}[H]
\centering
\caption{Training mode behavior by hints, feedback, and metrics}
\label{tab:training-mode-behavior-by-hints-feedback-and-metrics}
\small\renewcommand{\arraystretch}{1.2}\setlength{\tabcolsep}{4pt}
\begin{tabularx}{\textwidth}{|>{\raggedright\arraybackslash}X|>{\raggedright\arraybackslash}X|>{\raggedright\arraybackslash}X|>{\raggedright\arraybackslash}X|>{\raggedright\arraybackslash}X|}
\hline
Mode & Hints & Corrective Feedback & Metrics & Use \\
\hline
Pre-test & No & No & Yes & Baseline simulated tray-task proficiency. \\
\hline
Study cards & Yes & Yes & Optional & Prompt-before-reveal retrieval of names, aliases, families, features, and counts. \\
\hline
Quiz & Limited after answer & Yes & Yes & Retrieval practice, confidence calibration, and weak-item discovery. \\
\hline
Practice sort & Yes & Yes & Yes & Applied retrieval and deliberate practice on tray organization. \\
\hline
Post-test & No & No & Yes & Comparable same-day assessment after practice. \\
\hline
Delayed retention & No & No & Yes & Checks whether gains persist after 7 days (with falsification criteria). \\
\hline
\end{tabularx}
\end{table}
\subsection{Acceptance Criteria}
\label{sec:9-5-acceptance-criteria}
The training module is evidence-ready when it can show:
\begin{itemize}
  \item a learner completing a no-hints pre-test and post-test on comparable tray tasks;
  \item practice mode producing error-specific feedback;
  \item quiz mode recording correctness, time, and confidence or \texttt{Not sure};
  \item weak-item review repeating missed or uncertain instruments if time allows;
  \item metrics export comparing pre/post accuracy, duration, confidence, and error counts;
  \item a marker reliability check showing the app is not mostly scoring camera errors;
  \item expert or instructor review of instrument names, counts, distractors, lookalikes, and feedback;
  \item documentation that results are simulated training evidence, not clinical competency evidence.
\end{itemize}
\section{Evaluation And Validation Plan}
\label{sec:10--evaluation-and-validation-plan}
The evaluation plan uses a ladder of studies. Each level earns a different
claim, which prevents the project from overreaching.
\subsection{Evaluation Framework}
\label{sec:10-1-evaluation-framework}
The question of how to prove that people have learned requires an evaluation
framework that makes explicit what counts as evidence, for which claim, and at
what level of certainty. The training-evaluation and educational-assessment
literatures provide well-established structures for this.

Kirkpatrick's four-level model is the most widely used framework for training
evaluation: Level 1 (Reaction), Level 2 (Learning), Level 3 (Behavior), and
Level 4 (Results) \cite{kirkpatrick-evaluation-1994}. In this model, a pre/post
accuracy gain on a simulated tray task is Level 2 evidence: it shows that
knowledge or skill was acquired during the training session. The professor's
question about whether accuracy scores translate to real work corresponds to
Level 3 (Behavior transfer), which requires measuring performance in or near the
actual work context. Arthur et al.'s meta-analysis of organizational training
confirms this structure: they found that the mean effect size for learning
outcomes (Level 2) was larger than for behavior (Level 3) or results (Level 4),
and that evaluation designs must match the claim level
\cite{arthur-et-al-2003}. Salas et al. further argue that effective training
should be treated as a system that analyzes work, targets specific
competencies, provides practice and feedback, and evaluates transfer
\cite{salas-et-al-2012}. The current study ladder addresses Level 2 directly
and builds toward Level 3 through convergent evidence and expert review, but
does not claim Level 4 (reduced OR delays or hospital defect rates), which
requires an SPD-site study.

Within educational assessment, Messick's unified validity theory defines
construct validity as a single concept supported by six aspects: content,
substantive, structural, generalizability, external, and consequential
\cite{messick-validity-1989}. The Standards for Educational and Psychological
Testing adopt this unified view, treating validity as the degree to which
evidence and theory support the intended interpretations of test scores
\cite{aera-standards-2014}. Cook and Beckman operationalize the same framework
for health-professions education as five sources of validity evidence: content,
response process, internal structure, relations to other variables, and
consequences \cite{cook-beckman-validity-2006}. Kane's argument-based approach
adds that validity depends on the plausibility of an interpretive argument that
connects test scores to their intended use, with assumptions that must be
explicitly warranted at each inference step: scoring, generalization,
extrapolation, and implications \cite{kane-validity-2013}.

Applied to TrayGuard, these frameworks require the following evidence chain:
\begin{enumerate}
  \item \textbf{Content validity}: the tray module, names, aliases, counts,
  distractors, and lookalike pairs are plausible for a training tray. Addressed
  by expert review (Study 3) and instructor-verified module content.
  \item \textbf{Response process}: learners interact with the system as intended;
  confidence ratings and \texttt{Not sure} capture reflect genuine self-assessment
  rather than confusion or guessing. Addressed by usability observation and
  confidence-calibration analysis (Studies 1-2).
  \item \textbf{Internal structure}: error categories behave consistently: lookalike
  pairs produce misidentified errors, distractors produce extra errors, and
  repeated practice reduces specific error types. Addressed by pre/post error-type
  analysis (Studies 1-2).
  \item \textbf{Relations to other variables}: TrayGuard tray-sort scores correlate
  with an independent measure of the same knowledge. This is the weakest link
  in the current design. Addressed by a convergent photo-identification test
  added to Study 1.
  \item \textbf{Consequences}: the scores support the intended interpretation
  (improved simulated tray familiarity) and do not support unintended
  interpretations (clinical competence, SPD readiness). Addressed by scope
  boundaries and report labeling (Section~\ref{sec:11--risk-assessment-and-scope-boundaries}).
\end{enumerate}

A sixth issue spans the last two sources: \textbf{population generalizability}.
Study 1 tests complete novices (college students with no SPD experience), whose
pre-test accuracy is near zero. The target population is SPD workers, whose
baseline accuracy on familiar trays may already be high. A pre/post gain from
0 to 90 in novices does not, by itself, predict a gain from 90 to 90 in
experienced workers. This is a classical range-restriction and ceiling-effect
problem in training evaluation \cite{arthur-et-al-2003}. Kane's extrapolation
inference requires evidence that the learning mechanism transfers across
populations even when the score distribution shifts
\cite{kane-validity-2013}.

There are three responses to this objection within the paper's scope:
\begin{enumerate}
  \item \textbf{The Ofstead precedent shows that experienced workers are not at
  ceiling on novel or specialized tasks.} Ofstead et al.'s sterile-processing
  training pilot used certified SP professionals whose mean pre-test score was
  41\% — well below ceiling — on a borescope visual-inspection task that was
  new to them \cite{ofstead-et-al-2023}. SPD workers are experts on their
  current trays but will face new modules, updated count sheets, loaner trays,
  or specialized instrument sets. TrayGuard's measurable value for an
  experienced population is not absolute accuracy gain (which may be small)
  but error-type reduction (eliminating the specific errors they do make),
  confidence calibration (reducing high-confidence errors), standardization
  across shifts (aligning names and rules), and retention maintenance
  (preventing skill decay between exposure to low-frequency trays).
  \item \textbf{The dependent variable changes with the population.} For
  novices, the primary outcome is pre/post accuracy gain (baseline → learning).
  For experienced workers, the primary outcomes shift to: error-type reduction
  (fewer missing, wrong, misidentified items on known trays), speed improvement
  without accuracy loss, confidence calibration (high-confidence error rate),
  new-module learning rate (time to criterion on an unfamiliar tray), and
  retention over time (preventing decay on low-frequency trays). Study 4 (SPD
  pilot) must specify which dependent variable applies, rather than
  reusing the novice-study design unchanged \cite{kirkpatrick-evaluation-1994}.
  \item \textbf{The learning mechanism is population-independent even if the
  effect size is not.} Retrieval practice, immediate error-specific feedback,
  and spaced repetition improve performance across learner levels in
  health-professions education
  \cite{roediger-and-karpicke-2006, dunlosky-et-al-2013, larsen-et-al-2009}.
  The novice study establishes that the mechanism works in this task domain.
  The SPD pilot establishes the effect size for the target population. Arthur
  et al.'s meta-analysis confirms that training effect sizes vary by trainee
  population, which is why the study ladder separates the two populations
  rather than extrapolating effect sizes \cite{arthur-et-al-2003}.
\end{enumerate}
The study ladder (Table~\ref{tab:study-ladder-by-participant-type-and-claim-supported})
maps each study level to the evidence source it earns. No single study provides
all sources; the claim is supported by the accumulated pattern.

\subsection{Current Evidence Available}
\label{sec:10-2-current-evidence-available}
The current evidence package includes:
\begin{itemize}
  \item literature-backed problem definition for SPD tray reconstruction and assembly errors;
  \item support for simulation, retrieval practice, flashcards, feedback, tangible card interfaces, and board/card-based healthcare education;
  \item a bounded final product concept based on local tray modules and AprilTag instrument cards, with the evaluation framework above clarifying which claims each evidence source supports;
  \item a subsystem architecture, data contract, mode behavior, and scoring rubric;
  \item a user-study plan that addresses content validity, response process, internal structure, relations to other variables through convergent evidence, and retention through a committed delayed check.
\end{itemize}
\subsection{Formative Novice Study Plan}
\label{sec:10-3-formative-novice-study-plan}
The first feasible study is a formative novice study with students or other
participants who do not have SPD tray experience. The study should evaluate
whether the AprilTag-card simulation teaches a bounded local tray module, not
whether participants are ready for real SPD work.
Use one tray module with 8-12 required instruments, a distractor pool, and
parallel pre/post tray variants. A reasonable first sample is 8-12 participants
for classroom evidence; a larger follow-up can use paired statistics once the
prototype and protocol are stable.
Study sequence:
\begin{enumerate}
  \item Consent and prior-experience questionnaire.
  \item No-hints pre-test: assemble the simulated tray from AprilTag instrument cards using a tray list or count-sheet-style prompt.
  \item Retrieval-first study cards for instrument names, aliases, distinguishing features, and counts.
  \item Short quiz with confidence or \texttt{Not sure} capture.
  \item Practice tray sort with immediate error-specific feedback.
  \item Weak-item review for missed or uncertain instruments if time permits.
  \item No-hints post-test on a comparable tray variant.
  \item Short interview about confusing instruments, feedback usefulness, and perceived realism.
\end{enumerate}
Primary measures are pre/post tray accuracy, duration, missing/wrong/extra/
misidentified/wrong-count errors, confidence calibration, high-confidence
errors, low-confidence correct answers, and weak-item recovery. The study
supports the training claim if most participants improve accuracy without a
severe-error increase, maintain or improve completion time, and can explain at
least one corrected error after practice.

The primary value of the novice study is \textbf{instrument validation}, not
population-effect estimation. The study answers five prerequisite questions
that apply to any subsequent population:
\begin{enumerate}
  \item Do the error categories behave as expected (lookalike confusion vs.\
  missing vs.\ extra items), or does the scoring rubric produce noise?
  \item Does the convergent photo-identification test show a moderate-to-strong
  correlation with tray-sort accuracy ($r \ge 0.50$)? If not, the tray scores
  do not measure instrument knowledge in \emph{any} population.
  \item Is confidence calibration measurable (high-confidence error rate,
  \texttt{Not sure} frequency), or do participants all guess at the same rate?
  \item Does usability friction (scan failures, card confusion, UI navigation)
  dominate the learning signal, or can participants interact with the system
  as intended?
  \item Which distractors are effective, which feedback messages help, and
  which instruments are confusing? These content-refinement findings are
  transferable to an SPD module.
\end{enumerate}
If these fail in novices, they fail everywhere, because they are properties
of the measurement instrument, not of the participant population. An SPD
pilot run before the measurement tool is validated produces ambiguous results:
a null finding could mean the system does not help, or it could mean the
scoring rubric, error categories, or convergent measure are inadequate.
The novice study decouples these two failure modes.

A corollary is that the novice study does \textbf{not} establish the effect
size for SPD workers. The dependent variable shifts from absolute accuracy
gain (novices) to error-type reduction, confidence calibration, and retention
maintenance (workers) (Section~\ref{sec:10-1-evaluation-framework}). The
novice results warrant a worker study by showing the instrument functions;
they do not substitute for a worker study.

\subsubsection{Limitations And Counterarguments}
\label{sec:10-3-1-limitations-and-counterarguments}

The instrument-validation argument above has several limitations that
deserve explicit acknowledgment.

\paragraph{The measurement instrument includes the card interface, not just
the scoring rubric.} Validating error categories and convergent correlations
with AprilTag cards does not validate the same scores with real instruments.
The scoring rubric (error types, confidence calibration, accuracy) is
format-independent — a lookalike swap produces the same error code whether
the instruments are cards or steel. But the card-specific layers
(scan reliability, physical card handling, card-order confusion) are
validated separately by Study 0 (marker reliability) and usability
observation, not by the learning outcomes. The convergent photo-ID test
deliberately uses a different response format (paper-based name-and-belonging
identification) to show that the score construct transfers across formats.
A subsequent SPD pilot with real instruments or local photographs
revalidates the instrument for that medium.

\paragraph{Novices make obvious errors; the interesting errors workers make
may not appear.} A novice might miss 6 of 10 instruments; an SPD worker
might misidentify one lookalike pair. Both errors produce the same
item-level code (\texttt{missing} and \texttt{misidentified} respectively).
The category itself is format-agnostic and works at both granularities. What
cannot be tested with novices is whether the category set \emph{covers} the
error space of expert performance — whether workers make errors that fall
outside the current taxonomy (e.g., orientation errors, pairing errors,
assembly-order errors). That requires expert review (Study 3) to identify
missing categories and SPD pilot (Study 4) to test whether the taxonomy
captures real worker mistakes.

\paragraph{The convergent photo-ID test may floor for novices.} If
participants cannot identify any instruments in photographs after the
training session, the photo-ID test fails as a convergent measure not
because the tray-sort scores are invalid but because the photo-ID test is
too hard. Two design features address this risk. First, the test is
administered \emph{after} the full training sequence (study cards, quiz,
practice sort, post-test), by which point participants should know at least
the most visually distinctive instruments. Second, the test includes a
confidence rating per item and an explicit \texttt{Not sure} option, so
even chance-level name recall still produces a calibration signal. The
transfer probe (unseen instruments from a different module) provides a
convergent check: if participants perform at chance on unseen instruments
but above chance on trained ones, the photo-ID test is measuring
session-specific learning rather than general instrument familiarity.

\paragraph{Novice confidence ratings may be arbitrary.} Learners with no
prior task experience have no basis for calibrated self-assessment.
Uncalibrated confidence at pre-test is expected; the useful signal is
\emph{change} in calibration from pre-test to post-test. If confidence
ratings are uniformly low or uniformly high before and after, calibration
is not measurable and this evidence source fails. That failure is
informative: it means the training session does not improve self-awareness,
which is itself a meaningful finding about the system's effect on
metacognition.

\paragraph{Usability with novices does not predict usability with SPD
workers.} Novices may tolerate scan failures, card confusion, or UI
friction that an SPD worker would reject. The usability observation in
Study 1 is therefore a lower bound: if novices find the system unusable,
it is certainly unusable for workers. The converse does not hold — positive
novice usability does not guarantee worker acceptance — which is why the
study ladder reserves a separate usability assessment for the SPD pilot
(Study 4).
\subsection{Source Patterns For The Study Design}
\label{sec:10-4-source-patterns-for-the-study-design}
\begin{table}[H]
\centering
\caption{Source study patterns and TrayGuard translations}
\label{tab:source-study-patterns-and-trayguard-translations}
\small\renewcommand{\arraystretch}{1.2}\setlength{\tabcolsep}{4pt}
\begin{tabularx}{\textwidth}{|>{\raggedright\arraybackslash}X|>{\raggedright\arraybackslash}X|>{\raggedright\arraybackslash}X|}
\hline
Source & Study Pattern & TrayGuard Translation \\
\hline
NeuroVase used tangible cue cards, a tablet-based AR system, a structured curriculum, a controlled user study with 40 participants, a traditional-learning comparison, pre/post knowledge testing, SUS, Likert experience measures, and open feedback \cite{jahani-et-al-2026-neurovase}. & Use college students only as target-adjacent novices. Compare learning before and after the system, collect usability data, and route clinical transfer to later SPD validation. \\
\hline
Ofstead et al. used sterile-processing professionals, pre-testing, structured teaching, hands-on practice, confidence/satisfaction capture, workplace homework, and a 2-month booster/retention check \cite{ofstead-et-al-2023}. & Use the same pre/post/hands-on structure now, and treat a later SPD educator or SPD worker study as the stronger validation path. \\
\hline
Bell et al. randomized residents to post-tests at 0, 1, 3, 8, 21, or 55 days and found immediate gains could decay quickly; they recommended reinforcement after as little as 1 week \cite{bell-et-al-2008}. & Do not call same-day results retention. Add a 7-14 day delayed tray-sort check for a credible short-term retention claim. \\
\hline
Barrison et al. found electronic flashcards are widely used in health-professions education, while noting that development and delivery methods are less systematically studied \cite{barrison-et-al-2025}. & Use flashcards as a supported retrieval mechanism, but make the main outcome applied tray sorting rather than isolated card recall. \\
\hline
Ward et al. evaluated a serious board game for patient-safety education in two teaching hospitals using card-based discussion and learner feedback \cite{ward-et-al-2019}. & Treat physical cards as a legitimate educational abstraction when the claim is protocol learning or simulated decision-making, not patient outcome improvement. \\
\hline
\end{tabularx}
\end{table}
\subsection{Study Ladder}
\label{sec:10-5-study-ladder}
\begin{table}[H]
\centering
\caption{Study ladder by participant type and validity source}
\label{tab:study-ladder-by-participant-type-and-claim-supported}
\small\renewcommand{\arraystretch}{1.2}\setlength{\tabcolsep}{4pt}
\begin{tabularx}{\textwidth}{|>{\raggedright\arraybackslash}X|>{\raggedright\arraybackslash}X|>{\raggedright\arraybackslash}X|>{\raggedright\arraybackslash}X|}
\hline
Study & Participants & Validity source & Kirkpatrick level \\
\hline
Study 0: marker reliability & 1-2 project members, no learning claim & Measurement reliability (prerequisite) & -- \\
\hline
Study 1: novice simulated learning & 8-12 college students or classmates with no SPD experience & Response process, internal structure, relations to other variables (convergent) & Level 2 (Learning) \\
\hline
Study 2: delayed retention & Same Study 1 participants, 7 days later & Relations to other variables (temporal) & Level 2 (Learning, durability) \\
\hline
Study 3: expert review & 1 SPD educator, instructor, or knowledgeable reviewer & Content & Level 2 (Learning, content plausibility) \\
\hline
Study 4: future SPD pilot & SPD trainees or technicians & Consequential, extrapolation to target context & Level 3 (Behavior) \\
\hline
\end{tabularx}
\end{table}
The project should prioritize Studies 0-2 and, if possible, one expert review.
\subsection{Study 0: Marker Reliability Check}
\label{sec:10-6-study-0-marker-reliability-check}
Purpose: prevent app or camera failures from contaminating learner scores.
Setup:
\begin{itemize}
  \item Print all AprilTag cards used in the tray module.
  \item Use the same device, camera angle, lighting, tray boundary, and table setup planned for the user study.
  \item Test cards individually, in valid tray layouts, and in cluttered layouts with distractors.
\end{itemize}
Record total cards visible, cards detected, false positives, missed cards,
duplicate or unstable detections, and lighting or occlusion notes.
Acceptance target before user testing:
\begin{itemize}
  \item at least 95\% card detection in normal layouts;
  \item no systematic confusion between marker IDs;
  \item visible confirmation screen before final scoring so camera misses are not counted as learner errors.
\end{itemize}
\subsection{Study 1: Novice Simulated Learning}
\label{sec:10-7-study-1-novice-simulated-learning}
Recruit 8-12 college students or classmates with no formal SPD experience and
no prior exposure to the module. Collect optional background variables such as
medical, biology, lab, tool, or surgical-instrument familiarity. This matches
the NeuroVase-style proxy choice: students are not SPD workers, but they are
acceptable for a first novice learning and usability study if the claim is
bounded to simulated learning.
Use a within-subject pre/post design:
\begin{lstlisting}
intake -> pre-test -> training -> practice -> immediate post-test -> survey/interview
\end{lstlisting}
Target one 60 minute session per participant:
\begin{table}[H]
\centering
\caption{Study 1: novice session timeline}
\label{tab:study-1-novice-session-timeline}
\small\renewcommand{\arraystretch}{1.2}\setlength{\tabcolsep}{4pt}
\begin{tabularx}{\textwidth}{|>{\raggedright\arraybackslash}X|>{\raggedright\arraybackslash}X|>{\raggedright\arraybackslash}X|}
\hline
Segment & Time & Activity \\
\hline
Consent and intake & 3-5 min & Explain simulated study, collect background. \\
\hline
Device orientation & 2 min & Show how to scan and submit without teaching tray answers. \\
\hline
Pre-test tray sort & 5-8 min & No hints, no feedback, timed. \\
\hline
Retrieval cards & 10-12 min & Prompt-before-reveal study of required items, distractors, features, and counts. \\
\hline
Quiz & 5-8 min & Short recall/recognition quiz with confidence or \texttt{Not sure}. \\
\hline
Practice sort & 8-10 min & Same style task with immediate feedback. \\
\hline
Immediate post-test & 5-8 min & Parallel variant, no hints, no feedback, timed. \\
\hline
Convergent photo-ID test & 5-7 min & Paper-based photo identification of each required instrument plus 2 unseen transfer-probe instruments. \\
\hline
Survey and interview & 5-8 min & SUS or short 5-point ease/usefulness/confidence survey plus open questions. \\
\hline
\end{tabularx}
\end{table}
Primary learning measures:
\begin{itemize}
  \item tray accuracy: required items correctly selected with correct count;
  \item severe errors: missing, wrong, misidentified, extra, wrong-count;
  \item duration from task start to final submission;
  \item high-confidence errors;
  \item \texttt{Not sure} selections or low-confidence correct answers;
  \item weak-item recovery from quiz/practice to post-test.
\end{itemize}
Convergent evidence measures:
\begin{itemize}
  \item \textbf{Photo-identification test}: after the post-test tray sort, learners
  complete a short paper-based test showing high-resolution photos of each
  required instrument and 2-3 distractor instruments. For each photo, they write
  the instrument name, indicate whether it belongs in the tray, and rate their
  confidence. This test shares the knowledge domain (instrument identity, aliases,
  lookalike discrimination) with the tray sort while differing in response format
  (identification without physical card selection). The target correlation
  between post-test tray accuracy and photo-ID accuracy is $r \ge 0.50$, following
  the convergent-validity standard that two measures of the same construct should
  show moderate-to-strong association \cite{messick-validity-1989, cook-beckman-validity-2006, downing-validity-2003}.
  \item \textbf{Transfer probe}: the photo-ID test includes 2 instruments from a
  different tray module that the learner has never seen. If TrayGuard teaches
  general instrument-discrimination skill rather than module-specific
  memorization, learners who performed well on the trained tray should show
  better discrimination on unseen instruments. This is a weak test
  (novices have no basis to identify unseen instruments), but the pattern across
  participants is informative: even chance-level performance on the transfer
  probe, when combined with strong within-module gains, supports the bounded
  claim of module-specific learning.
\end{itemize}
Usability and experience measures:
\begin{itemize}
  \item completion rate;
  \item help requests;
  \item scan failures or rescan attempts;
  \item SUS or short 5-point ease/usefulness/confidence survey;
  \item open feedback on confusing cards, confusing instruments, and perceived realism.
\end{itemize}
Prior user studies support this design pattern. NeuroVase is the closest
interaction precedent because it tested a tangible cue-card and tablet-based AR
medical-learning system with 40 participants, a traditional-learning comparison,
pre/post knowledge testing, SUS, Likert experience measures, and open feedback
\cite{jahani-et-al-2026-neurovase}.
That supports TrayGuard's use of college students as target-adjacent novice
learners, provided the claim remains about simulated learning and usability
rather than target-worker competence.
Medical simulation app studies also commonly use novice learner cohorts when
the goal is early validity, usability, or learning evidence. Tulipan et al.
evaluated Touch Surgery's carpal-tunnel-release module with medical students,
orthopedic residents, and expert hand surgeons; medical students served as the
novice cohort, participants completed repeated simulation attempts, performance
was logged, and the novice group completed a Likert satisfaction/face-validity
questionnaire \cite{tulipan-et-al-2019}.
The study is relevant because TrayGuard similarly expects repeated attempts,
objective scoring, and user perception data to show whether the simulation is
usable and whether performance changes with practice.
Orthopedic VR training studies provide a second proxy-participant precedent:
Orland et al. studied first- and second-year medical students as novices in a
controlled orthopedic training task, comparing simulation-supported preparation
with conventional preparation and measuring downstream task performance
\cite{orland-et-al-2020}. This supports using
students when the target is an early, controlled, preclinical learning task,
while still reserving real workplace transfer for later target-user studies.
Finally, Ofstead et al. is the target-domain model rather than the proxy-user
model. Their sterile-processing training pilot used certified SP professionals,
pre/post testing, hands-on practice, confidence and satisfaction capture,
workplace homework, and a delayed booster
\cite{ofstead-et-al-2023}. TrayGuard's
student study should therefore be described as a first-stage novice simulation
study, with a later SPD educator review and SPD trainee pilot needed before
making target-worker claims.

A stronger later design would compare flashcard-only practice against
flashcards plus simulated tray sorting in a between-group design:
\begin{table}[H]
\centering
\caption{Between-group study design for isolating practice effect}
\label{tab:between-group-study-design-for-isolating-practice-effect}
\small\renewcommand{\arraystretch}{1.2}\setlength{\tabcolsep}{4pt}
\begin{tabularx}{\textwidth}{|>{\raggedright\arraybackslash}X|>{\raggedright\arraybackslash}X|>{\raggedright\arraybackslash}X|}
\hline
Group & Training condition & Post-test \\
\hline
A & Structured flashcards and quizzes only. & No-hints simulated tray-verification task. \\
\hline
B & Structured flashcards and quizzes plus simulated tray-verification practice with feedback. & Same no-hints simulated tray-verification task. \\
\hline
\end{tabularx}
\end{table}
The main outcome would be whether Group B performs better on the applied tray
task, measuring tray-verification accuracy, completion time,
missing/extra/wrong/misidentified/wrong-count errors, high-confidence errors,
low-confidence correct answers, and weak-item recovery. This design would
isolate the incremental effect of simulated practice beyond retrieval cards
alone. For the first project prototype, the within-subject pre/post design provides
simpler logistics and smaller sample requirements; the between-group comparison
is a natural follow-up once the prototype and protocol are stable.
\subsection{Study 2: Delayed Retention Check}
\label{sec:10-8-study-2-delayed-retention-check}
Same-day post-testing supports immediate learning only. It should not be called
retention \cite{bell-et-al-2008}. The delayed check distinguishes immediate
practice gain from retained learning. In Kane's validity argument framework, the
extrapolation inference from test score to real-world competence depends on
evidence that performance persists after a delay during which transient factors
(memorization of exact layouts, test familiarity, short-term motivational
effects) have decayed \cite{kane-validity-2013}. Without a delayed measure, the
interpretive argument supports only a same-day practice-effect claim.

The primary retention interval is \textbf{7 days} from the initial session. This
interval is a practical compromise: Bell et al. found that medical-knowledge
gains from a single intervention could decay measurably within 7 days
\cite{bell-et-al-2008}, and Larsen et al. showed that the testing-effect
advantage over restudy persisted at least 7 days in a medical-education
randomized trial \cite{larsen-et-al-2009}. Ofstead et al. used a longer
(2-month) interval with a booster session, but their intervention included
workplace homework between sessions \cite{ofstead-et-al-2023}. For a single-
session training tool without between-session homework, 7 days is the shortest
credible retention interval that separates durable learning from ephemeral
practice gain.

Procedure:
\begin{enumerate}
  \item Ask whether the participant studied the material since the first session.
  \item Run a no-hints delayed tray sort with a third parallel variant.
  \item Capture confidence or \texttt{Not sure}.
  \item Run the photo-identification test again to measure convergent-retention correlation.
  \item Ask what they remembered and what decayed.
\end{enumerate}

Falsification criteria. The retention claim is supported only if:
\begin{itemize}
  \item delayed accuracy retains at least half of the immediate pre/post gain
  (i.e., $\text{delayed} - \text{pre} \ge 0.5 \times (\text{post} - \text{pre})$);
  \item delayed accuracy does not return to within 5 percentage points of the
  pre-test baseline (i.e., the retention effect is not attributable to measurement
  noise or test-retest practice);
  \item at least 70\% of Study 1 participants complete Study 2.
\end{itemize}
If these criteria are not met, the project claim reverts to immediate learning
only and the retention gap is reported as a limitation that future work
(addressable by distributed practice, multiple training sessions, or a booster)
must close. This follows Kane's principle that the interpretive argument must
include explicit rebuttals: the conditions under which the claim would not be
supported \cite{kane-validity-2013}.

Optional extended intervals. If the 7-day check succeeds and schedule permits,
additional checks at 21-28 days or 2 months would strengthen the retention claim
further. The 2-month interval would match the Ofstead booster precedent
\cite{ofstead-et-al-2023}.
\subsection{Study 3: Expert Review}
\label{sec:10-9-study-3-expert-review}
Expert review addresses the main weakness of college-student testing: content
validity. The reviewer can be an SPD educator, SPD supervisor,
sterile-processing instructor, surgical technology instructor, or faculty
member with relevant instrument knowledge.
Materials to show:
\begin{itemize}
  \item tray list;
  \item instrument cards;
  \item distractor list;
  \item aliases and distinguishing features;
  \item scoring rubric;
  \item sample learner report.
\end{itemize}
Questions:
\begin{enumerate}
  \item Are the instrument names and aliases plausible?
  \item Are the required counts and distractors plausible for a training tray?
  \item Are the lookalike pairs educationally meaningful?
  \item Are any feedback messages misleading or unsafe?
  \item Would the weak-item report help target coaching?
  \item What would need to change before using this with SPD trainees?
\end{enumerate}
\subsection{Analysis And Acceptance Bar}
\label{sec:10-10-analysis-and-acceptance-bar}
For 8-12 participants, report descriptive paired results:
\begin{itemize}
  \item participant-level pre/post table;
  \item mean and median accuracy change;
  \item mean and median duration change;
  \item total error counts by type before and after;
  \item high-confidence errors before and after;
  \item weak-item examples;
  \item representative qualitative feedback.
\end{itemize}
Avoid strong inferential claims unless the sample is larger. If the sample is
large enough, add paired differences with confidence intervals.
The training claim is supported if most learners show:
\begin{itemize}
  \item higher post-test accuracy;
  \item equal or lower post-test duration without accuracy decline;
  \item fewer severe errors, especially missing and misidentified items;
  \item fewer high-confidence errors;
  \item fewer \texttt{Not sure} responses or low-confidence correct answers;
  \item evidence that missed or uncertain items improve after repeated retrieval;
  \item interview evidence that they can name what they learned or what remains hard.
\end{itemize}
Best project claim wording:
\begin{quote}
In a novice college-student sample, TrayGuard improved immediate simulated
tray-sorting performance after retrieval-centered practice. A delayed
7-day check tested short-term retention. A convergent photo-identification test
provided evidence that the tray-sort scores reflect instrument knowledge rather
than task-specific familiarity. Expert review was used to assess whether the
tray content and feedback were plausible for future SPD trainee validation.
These findings establish that the learning mechanism works in this task domain
for a naive population. They do not establish the effect size for experienced
SPD workers, whose baseline on familiar trays may be high and for whom the
dependent variable shifts from accuracy gain to error-type reduction,
confidence calibration, and retention maintenance (Section~10-1).
\end{quote}
\subsection{Validity Claims Summary}
\label{sec:10-11-validity-claims-summary}
The following table maps each project claim to the type of validity evidence
that supports it, following the Cook and Beckman five-source framework
\cite{cook-beckman-validity-2006} and Kane's argument-based approach
\cite{kane-validity-2013}. No single study provides all sources; the overall
claim rests on the accumulated pattern.
\begin{table}[H]
\centering
\caption{Validity claims mapped to evidence sources}
\label{tab:validity-claims-mapped-to-evidence-sources}
\small\renewcommand{\arraystretch}{1.2}\setlength{\tabcolsep}{4pt}
\begin{tabularx}{\textwidth}{|>{\raggedright\arraybackslash}X|>{\raggedright\arraybackslash}X|>{\raggedright\arraybackslash}X|>{\raggedright\arraybackslash}X|}
\hline
Claim & Validity source & Evidence & Would be falsified by \\
\hline
The tray module content is plausible for training. & Content & Study 3 expert review; instructor-verified module; literature grounding and real count-sheet fields (Section~\ref{sec:9-1-2-real-tray-list-examples}). & Expert finds names, counts, distractors, or feedback unrealistic; module uses content that contradicts published tray listings. \\
\hline
Learners interact with the system as intended; confidence reflects genuine self-assessment. & Response process & Study 1 usability observation; confidence-calibration analysis (high-confidence error rate, \texttt{Not sure} frequency); SUS scores; interview feedback about confusion points. & Frequent scan failures produce scores that reflect app error rather than learner knowledge; confidence ratings do not distinguish correct from incorrect items; SUS scores below 68. \\
\hline
Error categories behave consistently and decompose learning gains. & Internal structure & Study 1-2 pre/post error-type analysis; lookalike pairs produce proportionally more misidentified errors than non-lookalike items; practice reduces specific error categories rather than shifting errors across types. & Error types do not differ between lookalike and non-lookalike pairs; accuracy gain comes entirely from one error category (e.g., fewer missing items) while other categories remain unchanged. \\
\hline
Tray-sort scores measure instrument knowledge and not only task-specific familiarity. & Relations to other variables & Study 1 convergent photo-ID test; correlation between tray-sort accuracy and photo-ID accuracy; transfer probe on unseen instruments. & Photo-ID correlation $r < 0.30$; photo-ID scores show ceiling or floor effects that prevent discrimination; transfer-probe performance is negatively correlated with trained-tray performance. \\
\hline
Learning persists beyond the initial session. & Relations to other variables (temporal) & Study 2 delayed tray sort at 7 days; retention relative to pre-test baseline; comparison of delayed accuracy to immediate post-test accuracy. & Delayed accuracy returns to within 5 pp of pre-test; less than half of the immediate gain is retained; fewer than 70\% of Study 1 participants return. \\
\hline
The report supports instructor coaching without supporting overclaim. & Consequences & Scope boundaries (Section~\ref{sec:11--risk-assessment-and-scope-boundaries}); report labels frame results as simulated practice evidence; export fields include weak items and confidence calibration. & Report language uses pass/fail or clinical-readiness labels; results are used for employment decisions; learners report feeling surveilled rather than coached. \\
\hline
Learning gains in novices transfer to SPD workers (extrapolation). & Generalizability / extrapolation & Ofstead precedent: experienced SP workers were well below ceiling on novel tasks; mechanism (retrieval practice + feedback) is population-independent per the training literature; dependent variable shifts for workers. & Worker study fails to show improvement on any worker-relevant outcome; novice-effect-size extrapolation is cited without worker data. \\
\hline
\end{tabularx}
\end{table}
\subsection{Future SPD Workplace Pilot Plan}
\label{sec:10-11-future-spd-workplace-pilot-plan}
A later SPD pilot requires local count sheets, local photos or real teaching
instruments, SPD educator review, technician or trainee participants,
approval/privacy review, and a comparison condition such as current study
materials or preceptor-led orientation. It should measure whether the system
helps target users in a workplace-relevant training context. Only that later
pilot can support claims about SPD trainee usefulness, supervised transfer, or
operational training adoption.

The SPD pilot must change the dependent variable from the novice study. Where
novices are measured on absolute accuracy gain (pre 0 to post 90), experienced
workers are measured on: error-type reduction (fewer specific errors on known
trays), speed improvement without accuracy loss, confidence calibration
(high-confidence error rate), new-module learning rate (time to criterion on
an unfamiliar tray), and retention maintenance (preventing decay on
low-frequency trays). The pilot protocol must specify which of these applies
and the relevant minimum viable improvement before data collection
\cite{kirkpatrick-evaluation-1994, arthur-et-al-2003}.
\subsection{Validation Claims By Evidence Level}
\label{sec:10-12-validation-claims-by-evidence-level}
\begin{table}[H]
\centering
\caption{Validation claims by evidence level, validity source, and falsification criteria}
\label{tab:validation-claims-by-evidence-level-and-falsification-criteria}
\small\renewcommand{\arraystretch}{1.2}\setlength{\tabcolsep}{4pt}
\begin{tabularx}{\textwidth}{|>{\raggedright\arraybackslash}X|>{\raggedright\arraybackslash}X|>{\raggedright\arraybackslash}X|>{\raggedright\arraybackslash}X|>{\raggedright\arraybackslash}X|}
\hline
Claim & Validity source & Current Support & Required Validation & What Would Falsify It \\
\hline
Marker cards are a usable proxy for simulated tray choices. & Measurement reliability & NeuroVase, QR/marker education precedent, and marker reliability testing. & Study 0 and participant workflow observation. & Frequent detection failures, answer leakage, or user confusion about card meaning. \\
\hline
Retrieval cards and quizzes improve immediate tray-task performance. & Internal structure & Retrieval-practice and flashcard literature. & Study 1 pre/post accuracy and error results. & No accuracy gain, more severe errors, or gains only from answer leakage. \\
\hline
Tray-sort scores reflect instrument knowledge. & Relations to other variables & Literature on simulation-based knowledge assessment. & Study 1 convergent photo-ID correlation $r \ge 0.50$. & Photo-ID correlation $r < 0.30$ or ceiling/floor effects. \\
\hline
Learning persists beyond the session. & Relations to other variables (temporal) & Retention literature and Ofstead booster precedent. & Study 2 7-day delayed tray sort with falsification criteria. & Delayed performance returns to within 5 pp of baseline; less than half of gain retained. \\
\hline
Content is plausible for SPD training. & Content & Local count-sheet and sterile-processing education literature. & Study 3 expert review. & Expert finds names, counts, distractors, or feedback unrealistic. \\
\hline
Novice gains extrapolate to SPD workers. & Generalizability / extrapolation & Mechanism evidence (retrieval practice works across populations) and Ofstead precedent. & Future SPD pilot with worker-appropriate dependent variable. & Worker study shows no improvement on any worker-relevant outcome. \\
\hline
TrayGuard helps SPD trainees. & Consequential & Not proven by the current student study. & Future SPD pilot. & SPD trainees reject workflow or fail to improve on workplace-relevant tasks. \\
\hline
\end{tabularx}
\end{table}
\subsection{Design Outcome Versus Requirements}
\label{sec:10-13-design-outcome-versus-requirements}
\begin{table}[H]
\centering
\caption{Design outcome versus requirements status}
\label{tab:design-outcome-versus-requirements-status}
\small\renewcommand{\arraystretch}{1.2}\setlength{\tabcolsep}{4pt}
\begin{tabularx}{\textwidth}{|>{\raggedright\arraybackslash}X|>{\raggedright\arraybackslash}X|}
\hline
Requirement Area & Current Status \\
\hline
Local tray module structure & Ready for implementation. \\
\hline
Retrieval-first study and quiz modes & Ready for implementation. \\
\hline
AprilTag-card tray sorting & Ready for prototype build and marker reliability testing. \\
\hline
Error-specific scoring & Defined; ready for implementation. \\
\hline
Pre/post novice study with convergent photo-ID & Protocol-ready. \\
\hline
Delayed retention with falsification criteria & Protocol-ready. Committed 7-day check. \\
\hline
Expert content review & Ready if an instructor or SPD educator is available. \\
\hline
SPD trainee validation & Future clinical/workplace validation required. \\
\hline
Deployment-grade real-instrument CV & Future support-layer validation required. \\
\hline
\end{tabularx}
\end{table}
\section{Risk Assessment And Scope Boundaries}
\label{sec:11--risk-assessment-and-scope-boundaries}
The main risk is not that the design is impossible. The main risk is
overclaiming. TrayGuard is viable as a simulated learning prototype, but the project
must separate immediate novice learning, delayed retention, expert content
review, future SPD transfer, and future real-instrument CV.
\subsection{Formal Risk Register}
\label{sec:11-1-formal-risk-register}
\begin{table}[H]
\centering
\caption{Formal risk register with causes, consequences, and mitigations}
\label{tab:formal-risk-register-with-causes-consequences-and-mitigations}
\small\renewcommand{\arraystretch}{1.15}\setlength{\tabcolsep}{3pt}
\begin{tabularx}{\textwidth}{|>{\raggedright\arraybackslash}p{2.2cm}|>{\raggedright\arraybackslash}X|>{\raggedright\arraybackslash}p{2.5cm}|>{\raggedright\arraybackslash}X|>{\raggedright\arraybackslash}p{2.3cm}|}
\hline
Risk & Cause & Consequence & Mitigation & Residual Boundary \\
\hline
Student proxy overclaim & Participants are not SPD workers. & Reviewers reject relevance to target users. & State that students test novice simulated learning; add expert review; reserve SPD claims for future pilot. & No SPD worker effectiveness claim. \\
\hline
Same-day retention mistake & Immediate post-test follows practice. & Report overstates durability. & Add committed 7-day delayed check with falsification criteria; otherwise call it immediate learning only. & Retention claim only after delayed test. \\
\hline
Card physical realism & Cards do not reproduce weight, scale, tactile handling, or real visual lookalikes. & Transfer to real instruments remains unknown. & Use cards as selection/count proxy; add local photos or real teaching instruments later. & No clinical handling competency claim. \\
\hline
Marker detection errors & Lighting, occlusion, angle, or overlap causes missed detections. & App errors look like learner errors. & Run Study 0; log raw detections; show confirmation before scoring. & Marker reliability must be reported. \\
\hline
Answer leakage from cards & Labels or marker IDs reveal the answer. & Measured gain is not learning. & Use neutral assessment fronts; randomize marker IDs; hide answer text in assessment. & Study and assessment cards must be separated. \\
\hline
Content invalidity & Student-built module may use wrong names, counts, or distractors. & Results do not reflect plausible SPD training. & Require expert review of module content and feedback. & Content validity limited until SPD review. \\
\hline
Feedback contaminates assessment & Practice answers leak into post-test. & Pre/post comparison inflated. & Use parallel variants; suppress hints and feedback in assessment modes. & Cannot claim broad generalization from one tray. \\
\hline
Usability friction dominates learning & Scanning or workflow confusion slows participants. & Poor scores reflect interface friction. & Track help requests, scan failures, SUS/ease ratings, and qualitative feedback. & Usability results must be reported separately. \\
\hline
\end{tabularx}
\end{table}
\subsection{Unknowns And Concerns}
\label{sec:11-2-unknowns-and-concerns}
The current design is viable as a low-cost formative learning prototype, but
several concerns remain unresolved:
\begin{table}[H]
\centering
\caption{Unknowns and concerns with strengthening actions}
\label{tab:unknowns-and-concerns-with-strengthening-actions}
\small\renewcommand{\arraystretch}{1.2}\setlength{\tabcolsep}{4pt}
\begin{tabularx}{\textwidth}{|>{\raggedright\arraybackslash}X|>{\raggedright\arraybackslash}X|>{\raggedright\arraybackslash}X|}
\hline
Concern & Why it matters & Strengthening action \\
\hline
Card-to-instrument transfer & AprilTag cards preserve selection and counting decisions, but they do not reproduce weight, scale, handling, tactile inspection, or real visual lookalikes. & Treat the first study as simulated learning only; add an SPD educator review and a later task using local photos or real teaching instruments. \\
\hline
Flashcard overreach & Flashcards support recall, but tray work requires applying local rules under distractor pressure. & Keep flashcards as preparation, then make the main outcome a no-hints tray sort with missing, extra, wrong, misidentified, and wrong-count scoring. \\
\hline
Marker identity leakage & If cards visibly encode names or IDs, learners may match labels instead of learning instruments. & Use neutral card fronts during assessment, randomize marker IDs, hide answer text, and separate study cards from assessment cards. \\
\hline
Content validity & A student-built tray module may contain wrong names, unrealistic distractors, or nonlocal quantities. & Require instructor or SPD educator verification of every instrument, alias, count, and distractor before using results as evidence. \\
\hline
Assessment equivalence & Pre/post gains are weak evidence if the post-test is easier or repeats the same exact layout. & Build parallel pre/post tray variants with matched item families, counts, distractors, and difficulty. \\
\hline
Short-term learning only & A same-day post-test can show practice effects rather than retention. & Add a delayed retention check or booster if schedule permits. \\
\hline
Novice participant limits & Student participants can test learnability, but they do not represent SPD technicians under workflow pressure. & Report student results as novice simulated-learning evidence and reserve adoption/clinical-transfer claims for an SPD pilot. \\
\hline
Scanning reliability & Camera angle, lighting, occlusion, and card overlap may produce app errors that look like learner errors. & Log raw detections, show a confirmation state before scoring, and run a marker-detection bench test before the user study. \\
\hline
Confidence calibration & Confidence ratings can become noise if captured too often or too vaguely. & Capture simple confidence or \texttt{Not sure} at assessment moments and analyze high-confidence errors separately. \\
\hline
Instructor usefulness & Better learner scores do not automatically mean the output helps educators. & Add a short instructor-facing report and ask an educator whether the weak-item summary would change coaching. \\
\hline
\end{tabularx}
\end{table}
These holes do not invalidate the design. They define the claim boundary: the
prototype can show whether a marker-card training loop improves novice
performance on a simulated local tray task. It cannot yet show real SPD
competence, reduced OR delays, or deployment-ready instrument recognition.
\subsection{Scope Boundaries}
\label{sec:11-3-scope-boundaries}
Allowed after Study 1:
\begin{itemize}
  \item novices improved on a simulated tray task;
  \item the interface was usable for first-time learners;
  \item retrieval cards plus practice sorting produced measurable immediate gains.
\end{itemize}
Allowed after Study 2:
\begin{itemize}
  \item gains persisted after 7 days (meeting the falsification criteria in Section~\ref{sec:10-8-study-2-delayed-retention-check});
  \item the system showed measurable short-term retention in a simulated tray task, distinguishable from same-day practice effects.
\end{itemize}
Allowed after Study 3:
\begin{itemize}
  \item an expert reviewer found the module plausible for future SPD validation;
  \item content concerns were identified and revised before broader testing.
\end{itemize}
Not allowed yet:
\begin{itemize}
  \item SPD technicians will improve;
  \item the system reduces tray errors in hospitals;
  \item the system reduces OR delays;
  \item the system certifies competency;
  \item the app recognizes real surgical instruments;
  \item same-day gains demonstrate retention.
\end{itemize}
\section{Broader Impacts And Ethics}
\label{sec:12--broader-impacts-and-ethics}
TrayGuard's positive impact is a more accessible way to practice local tray
knowledge before a learner consumes scarce supervised training time. If it
works, learners get repeatable practice, instructors get clearer weak-item
evidence, and future teams get a safer bridge between training content and
camera-supported tray review.
The ethical risk is that a useful training report could be misused. A learner
score is not a certification result, employment screen, disciplinary record, or
permission to perform unsupervised sterile-processing work. The report should
therefore label outputs as simulated practice evidence and should avoid
language such as pass, fail, competent, or cleared unless a formal program adds
its own supervised competency process.
Using the principles of beneficence, nonmaleficence, autonomy, and justice:
\begin{table}[H]
\centering
\caption{Bioethics principles applied to TrayGuard}
\label{tab:bioethics-principles-applied-to-trayguard}
\small\renewcommand{\arraystretch}{1.2}\setlength{\tabcolsep}{4pt}
\begin{tabularx}{\textwidth}{|>{\raggedright\arraybackslash}X|>{\raggedright\arraybackslash}X|}
\hline
Principle & TrayGuard implication \\
\hline
Beneficence & The system should help learners improve and help educators target coaching, especially on missing, wrong, misidentified, and high-confidence errors. \\
\hline
Nonmaleficence & The system should not imply clinical readiness, replace hands-on supervision, or hide uncertainty behind a single score. \\
\hline
Autonomy & Learners and instructors should understand what is being recorded, why it is being recorded, and how the report will be used. \\
\hline
Justice & Access should not depend on having a perfect local instrument library, expensive hardware, or one preferred training site. Modules should be updateable and printable. \\
\hline
\end{tabularx}
\end{table}
Negative externalities and controls:
\begin{table}[H]
\centering
\caption{Negative externalities and ethical controls}
\label{tab:negative-externalities-and-ethical-controls}
\small\renewcommand{\arraystretch}{1.2}\setlength{\tabcolsep}{4pt}
\begin{tabularx}{\textwidth}{|>{\raggedright\arraybackslash}X|>{\raggedright\arraybackslash}X|>{\raggedright\arraybackslash}X|}
\hline
Risk & Ethical concern & Control \\
\hline
Punitive learner metrics & Scores could become surveillance or discipline instead of coaching. & Use hashed IDs for studies; report weak items and practice progress; prohibit employment or competency decisions from prototype data. \\
\hline
Overreliance on simulated scores & A high score could be mistaken for real SPD readiness. & Label all results as simulated; require supervised hands-on validation for clinical claims. \\
\hline
Reduced supervised practice & Programs might use the tool to cut preceptor time. & Frame TrayGuard as preparation for supervised practice, not replacement. \\
\hline
Unequal local content & Better-resourced sites may build richer modules. & Use open file formats, printable cards, and a minimal module template. \\
\hline
Instrument-library bias & Learners may only practice instruments represented in the module. & Version modules, list coverage gaps, and require local review before use. \\
\hline
Answer leakage & Labels or marker IDs could inflate learning results. & Use neutral assessment cards and hide names/counts in assessment modes. \\
\hline
Printing and material waste & Card iterations create paper/plastic waste. & Use reusable sleeves, print small modules, and revise digitally before reprinting. \\
\hline
Privacy creep & Future reports could collect names, staff IDs, or workplace behavior. & Keep patient data out of scope; collect only fields needed for learning analysis. \\
\hline
\end{tabularx}
\end{table}
These controls are part of the design, not optional documentation. They protect
the stakeholder value of TrayGuard by keeping the training tool honest about
what it can and cannot prove.
\section{Project Plan And Future Work}
\label{sec:13--project-plan-and-future-work}
This FDR is now the handoff plan. The next team should build only the minimum
evidence loop needed to run the study, then expand after the first results.
\subsection{Project Plan Postmortem}
\label{sec:13-1-project-plan-postmortem}
The original project direction emphasized computer vision support for tray
checking. That direction remained technically plausible, but the validation
burden was too large for the project schedule: real instrument recognition would
require local instrument data, site-specific tray rules, cross-manufacturer
generalization testing, camera reliability, clinical workflow review, and
adoption evidence. The project therefore pivoted to a training-centered system
with a stronger near-term evidence path.
The pivot reduced risk by changing the main proof obligation from
deployment-grade recognition to measurable novice learning. The remaining
engineering work is still concrete: build a local tray module, use AprilTag
cards to make the simulated physical task observable, score tray attempts, and
run a small pre/post study with claim boundaries.
\subsection{Future Work Roadmap}
\label{sec:13-2-future-work-roadmap}
\begin{table}[H]
\centering
\caption{Future work roadmap phases}
\label{tab:future-work-roadmap-phases}
\small\renewcommand{\arraystretch}{1.2}\setlength{\tabcolsep}{4pt}
\begin{tabularx}{\textwidth}{|>{\raggedright\arraybackslash}X|>{\raggedright\arraybackslash}X|>{\raggedright\arraybackslash}X|}
\hline
Phase & Goal & Exit Criteria \\
\hline
1. Minimal prototype & Build the end-to-end local tray loop. & Load module, scan cards, run pre-test, show study cards, run quiz, run practice sort, run post-test, export metrics. \\
\hline
2. Verified module & Make the first tray content defensible. & 8-12 required instruments, distractors, lookalikes, neutral assessment cards, matched variants, and expert/instructor review. \\
\hline
3. Marker reliability & Prove the app can observe card choices. & At least 95\% detection in normal layouts, no systematic ID confusion, confirmation screen before scoring. \\
\hline
4. Novice study & Test immediate simulated learning. & 8-12 college-student participants, pre/post results, usability notes, and bounded claim language. \\
\hline
5. Retention check & Test whether learning persists. & 7-day delayed no-hints tray sort with third matched variant, plus falsification criteria (Section~\ref{sec:10-8-study-2-delayed-retention-check}). \\
\hline
6. Report refinement & Turn study output into project evidence. & Participant-level results, aggregate pre/post plots, error-category table, high-confidence errors, and qualitative themes. \\
\hline
7. SPD educator review & Improve content validity. & Reviewer feedback on names, counts, distractors, lookalikes, scoring, and weak-item report. \\
\hline
8. Future SPD pilot & Test target-user relevance. & SPD trainee or technician participants, local count sheet/photos, approval/privacy review, and comparison condition. \\
\hline
9. Optional CV support & Reintroduce real-instrument recognition as support, not the core proof. & Detector tested on local photos with generalization and failure-mode reporting. \\
\hline
\end{tabularx}
\end{table}
\subsection{GitLab Epic / Issue Map}
\label{sec:13-3-gitlab-epic-issue-map}
\begin{table}[H]
\centering
\caption{Epic and issue map for continuation}
\label{tab:epic-and-issue-map-for-continuation}
\small\renewcommand{\arraystretch}{1.2}\setlength{\tabcolsep}{4pt}
\begin{tabularx}{\textwidth}{|>{\raggedright\arraybackslash}X|>{\raggedright\arraybackslash}X|}
\hline
Epic & Issues \\
\hline
Training Module & Define first tray module; add instrument cards with aliases and features; add tray quantities and distractors; build study-card view; build identification quiz. \\
\hline
Scoring And Feedback & Record quiz correctness, duration, and confidence; build simulated tray sorting; implement missing/extra/wrong/misidentified/wrong-count scoring; add practice feedback. \\
\hline
Assessment & Build pre-test mode; build post-test mode; suppress hints and feedback during assessment; add delayed-retention mode if schedule allows. \\
\hline
Evidence And Reporting & Export learner metrics; generate pre/post summary; report confidence calibration, high-confidence errors, and weak-item recovery. \\
\hline
Pilot Evidence & Write pilot walkthrough script; run marker reliability check; run novice study; record confusion points; add limitation note about non-SPD participants. \\
\hline
Content Validation & Prepare expert-review packet; collect reviewer feedback; revise names, counts, distractors, and feedback language. \\
\hline
Presentation & Lock final demo narrative around simulated learning, retention boundary, and future SPD validation. \\
\hline
\end{tabularx}
\end{table}
\subsection{Completed Work}
\label{sec:13-4-completed-work}

The following subsystems have been implemented:

\begin{itemize}
  \item \textbf{Training web app.} A FastAPI backend serves a single-page browser application with seven-step learner flow: intake, pre-test, study cards, quiz, practice sort, post-test, and summary. Camera capture through \texttt{getUserMedia} streams frames to a server-side ArUco detection endpoint. Manual quantity entry serves as a fallback when camera access is unavailable or unreliable \texttt{src/trayguard/web/app.py}.
  \item \textbf{Tray module data model.} Eight file-backed tray modules define required instruments, distractors, lookalike pairs, assessment variants, and ArUco marker mappings. The primary module, \texttt{basic\_general\_tray\_v1}, contains 10 required instrument types (15 total required units), 5 distractors, 6 lookalike pairs, and 2 assessment variants with matched difficulty. Five instrument catalogs covering 400+ classes provide authoring reference material \texttt{config/tray\_modules/}.
  \item \textbf{Card generation and printing.} A CLI command generates printable ArUco marker PNGs, a marker-grid index sheet for camera scanning, and double-sided study flashcards with instrument images, names, aliases, distinguishing features, and tray-specific quantities. An image-processing pipeline handles EXIF rotation, content-aware cropping, and automatic marker placement \texttt{src/trayguard/learning/cards.py}.
  \item \textbf{Scoring engine.} A scoring function classifies each item into correct, missing, extra, wrong, misidentified (lookalike substitution), or wrong-count categories. Error-point aggregation maps to an accuracy score and required-recall metric. Confidence ratings (1-5 scale and a \texttt{Not sure} checkbox) are captured per item and flagged as high-confidence errors or low-confidence correct answers \texttt{src/trayguard/learning/scoring.py}.
  \item \textbf{Export and run storage.} Each completed run produces a JSON payload plus CSV files at the attempt level and item level. Learner identifiers are hashed (SHA-256 truncated to 16 hex characters). The export schema includes all error categories, confidence metadata, and module version fields \texttt{src/trayguard/learning/export.py}; \texttt{src/trayguard/learning/run\_store.py}.
   \item \textbf{CV data collection and annotation pipeline.} A webcam capture tool with bounding-box annotation UI supports rapid multi-session data collection with per-class image output, YOLO label export, and session save/load \texttt{src/trayguard/data\_collection/}.
   \item \textbf{YOLO training infrastructure.} A Hydra-configured training pipeline supports model, dataset, and hyperparameter overrides across YOLO11n/s/m/l variants and RT-DETR-l. Training runs are deterministic and reproducible through a fixed seed \texttt{src/trayguard/training/}.
   \item \textbf{CV experiments: lighting robustness.} Seven staged brightness experiments (800--80 lumens) on matte, reflective, separated, and overlay layouts were trained and validated. Per-stage precision, recall, mAP50, mAP50-95, and per-image count-error metrics were collected across approximately 11 brightness levels \texttt{scripts/validate\_brightness\_experiments.py}.
   \item \textbf{CV experiments: shape similarity.} Six instrument-like classes in three pair families (straight versus curved scissors, straight versus curved clamps, narrow versus broad tip) were rendered in Blender and photographed as real proxies. Pairwise confusion matrices, high-confidence wrong-pair rates, and open-set unknown-distractor false positive rates were computed from synthetic-seen, synthetic-heldout, and real transfer evaluations \texttt{scripts/validate\_shape\_similarity\_experiments.py}.
   \item \textbf{CV experiments: synthetic-to-real transfer.} A real-proxy dataset was used to compare mAP50, wrong-similar-class rates, and unknown false positive rates between training from scratch and fine-tuning from a synthetic pretrained checkpoint \texttt{config/shape\_similarity/default\_experiment.json}.
   \item \textbf{Multi-variant model benchmark.} YOLO11n/s/m/l at 640 and 960 image sizes and RT-DETR-l at 640 were benchmarked for precision, recall, mAP, and latency \texttt{src/trayguard/training/benchmark.py}.
   \item \textbf{Legacy Streamlit tray-check demo.} A live YOLO tray-check interface demonstrates camera capture, detection, and tray-status comparison for proof-of-concept review \texttt{src/trayguard/app/streamlit\_app.py}.
   \item \textbf{Tests.} Unit tests verify module loading, scoring accuracy (full-credit and lookalike-substitution cases), CSV/JSON export file completeness, and ArUco marker generation and detection \texttt{tests/test\_learning.py}.
\end{itemize}
\subsection{Budget}
\label{sec:13-5-budget}
\begin{table}[H]
\centering
\caption{Project budget by category}
\label{tab:project-budget-by-category}
\small\renewcommand{\arraystretch}{1.2}\setlength{\tabcolsep}{4pt}
\begin{tabularx}{\textwidth}{|>{\raggedright\arraybackslash}X|>{\raggedright\arraybackslash}X|>{\raggedright\arraybackslash}X|}
\hline
Category & Expected Cost & Notes \\
\hline
Software prototype & \$0 direct cost & Uses existing local app stack and open-source AprilTag support where available. \\
\hline
Physical cards & \$5-25 & Printer paper/cardstock, sleeves, tape, or foam backing. \\
\hline
Tray surface & \$0-30 & Existing table, poster board, foam board, or simple marked tray boundary. \\
\hline
Camera/device & \$0 if using existing laptop/tablet/phone & Optional stand improves reliability but is not required for the first study. \\
\hline
Participant cost & \$0-100 & Small snacks/gift cards if allowed; class demo can be unpaid. \\
\hline
Expert review & \$0-200 & Likely unpaid instructor feedback; paid SPD educator review would be stronger. \\
\hline
Future SPD pilot & TBD & Depends on site approval, staff time, privacy review, and local instrument access. \\
\hline
\end{tabularx}
\end{table}
The main cost is labor: module authoring, prototype implementation, marker
testing, study moderation, analysis, and report writing.
\section{Discussion And Final Recommendations}
\label{sec:14--discussion-and-final-recommendations}
This report supports a clear next step: build and test the software-first training
loop before investing more effort in deployment-grade tray automation. The
problem evidence is strong enough to justify the target task, and the learning
evidence is strong enough to justify a simulated training prototype. The
current evidence is not strong enough to claim that TrayGuard reduces hospital
tray errors, improves SPD technician performance, or recognizes real surgical
instruments in clinical conditions.

The evaluation plan now follows established frameworks from the training-
evaluation and educational-assessment literatures. Kirkpatrick's four-level
model frames the project's reach: the current studies target Level 2 (Learning),
with internal-structure evidence from error-type decomposition, convergent
evidence from an independent photo-identification test, and temporal evidence
from a committed 7-day retention check. Level 3 (Behavior transfer) and Level 4
(Results) remain future work requiring SPD-site participants and operational
outcome data \cite{kirkpatrick-evaluation-1994}. The five-source validity
framework from Cook and Beckman, grounded in Messick's unified validity theory
and Kane's argument-based approach, provides the evidentiary structure for each
claim level, with explicit falsification criteria for the retention and
convergent-validity claims \cite{messick-validity-1989, kane-validity-2013,
cook-beckman-validity-2006, aera-standards-2014}.

A further validity concern is population generalizability. The novice study
tests complete novices (college students). Experienced SPD workers may already
perform at or near ceiling on their current trays. A pre/post gain from 0 to 90
in novices does not predict a gain from 90 to 90 in workers. The evaluation
design addresses this by separating the dependent variable by population:
novices are measured on absolute accuracy gain; workers would be measured on
error-type reduction, confidence calibration, new-module learning rate, and
retention maintenance (Section~\ref{sec:10-1-evaluation-framework}). The
Ofstead precedent shows that even certified SP workers have substantial room
for improvement on novel tasks (41\% pre-test in their study), supporting the
claim that TrayGuard's value for an experienced population lies in targeted
error reduction rather than general accuracy gain. The novice studies
establish that the learning mechanism works in this task domain; the SPD pilot
is required to establish the effect size for the target population.

Stakeholders can conclude the following:
\begin{itemize}
  \item Tray reconstruction is a defensible focus because many readiness failures involve assembly, visualization, identification, missing items, wrong items, extra items, and local count-sheet use.
  \item The training pivot is justified because local tray knowledge and supervised practice are real bottlenecks, while deployment-grade CV has a much larger validation burden.
  \item A marker-card prototype can support a bounded novice-learning study if the cards are neutral, detection is reliable, assessment variants are comparable, and results are labeled as simulated.
  \item The project has reduced concept risk, requirement risk, scoring risk, and study design risk by defining the module contract, performance targets, error categories, study ladder, convergent validity target, and claim boundaries.
  \item The evaluation design addresses content validity (expert review), internal structure (error-type analysis), relations to other variables (convergent photo-ID test), and temporal persistence (committed 7-day retention check with falsification criteria).
\end{itemize}
Stakeholders cannot yet conclude that TrayGuard:
\begin{itemize}
  \item certifies sterile-processing competence;
  \item replaces hands-on hours, preceptor observation, or local sign-off;
  \item reduces OR delays or hospital tray-defect rates;
  \item works for SPD technicians in a real workplace;
  \item recognizes real instruments across manufacturers, lighting, occlusion, and tray layouts;
  \item produces the same effect size for experienced workers as for novices (the dependent variable and expected magnitude differ by population);
  \item supports retention unless the 7-day delayed check meets the falsification criteria specified in Section~\ref{sec:10-8-study-2-delayed-retention-check}.
\end{itemize}
The next team should therefore implement the minimum evidence loop:
\begin{lstlisting}
load module -> pre-test -> study cards -> quiz -> practice sort -> post-test -> convergent photo-ID test -> export report
\end{lstlisting}
After that loop works, the team should run the marker reliability check, then a
small novice pre/post study with the convergent photo-ID measure, then the
committed 7-day retention check, then an expert review of names, counts,
distractors, lookalikes, and feedback. Only after those results should the
project expand toward real teaching instruments, local photos, SPD trainees, or
real-instrument CV.
Final recommendation:
\begin{quote}
Treat TrayGuard as a local tray learning and assessment platform first. Use
the project prototype to prove that novices can improve on a simulated local tray
task, demonstrate that the gains reflect instrument knowledge (convergent
photo-ID evidence), persist for at least 7 days (retention check), and are
grounded in plausible content (expert review), then use those results to decide
whether a future SPD pilot or camera-supported tray review system is worth the
additional validation cost.
\end{quote}
\section{References}
\label{sec:references}
The working source list is consolidated in
[bibliography.md](bibliography.md). Citations in this report should continue to
link to that file so a future team can update sources without rebuilding the
paper's structure.
\appendix
\section{Research Support Audit}
\label{sec:appendix-a-research-support-audit}

\begin{table}[H]
\centering
\caption{Research support audit: design choice traceability}
\label{tab:research-support-audit-design-choice-traceability}
\small\renewcommand{\arraystretch}{1.2}\setlength{\tabcolsep}{4pt}
\begin{tabularx}{\textwidth}{|>{\raggedright\arraybackslash}X|>{\raggedright\arraybackslash}X|>{\raggedright\arraybackslash}X|}
\hline
Design Choice & Research Support & Defensible Stance \\
\hline
Time-to-competency as the north star & SPD certification and training evidence shows that entry into practice requires hands-on hours, structured training, and preceptor support \cite{hspa-crcst-accessed-2026}; \cite{chobin-2010}; \cite{aorn-staffing-shortage-2025}. & Measure simulated improvement in local tray familiarity and report the longitudinal onboarding question as future field validation. \\
\hline
Retrieval-centered learning loop & Simulation, retrieval practice, spacing, feedback, and pre/post testing are supported by health-professions and learning-science evidence, with direct SP training precedent in Ofstead et al. and digital-practice support from spaced education and electronic-flashcard research \cite{ofstead-et-al-2023}; \cite{cook-et-al-2011}; \cite{roediger-and-karpicke-2006}; \cite{larsen-et-al-2009}; \cite{martinengo-et-al-2024}; \cite{barrison-et-al-2025}; \cite{hattie-and-timperley-2007}. & Build pre-test, retrieval-first cards, quizzes, applied tray practice with feedback, weak-item review, post-test, and metrics export as one module. \\
\hline
Local tray authoring & Instrument families are reusable, but actual tray requirements are local: count sheets specify tray contents, quantities, sizes, and catalog/reference numbers; tray optimization depends on procedure, surgeon preference, usage likelihood, and stock decisions; specialty and loaner trays can be vendor-specific \cite{nadeau-2024}; \cite{dos-santos-et-al-2021}; \cite{ahmadi-et-al-2023}. & Start with file-backed instructor-verified modules and record local names, aliases, quantities, photos, and version history. \\
\hline
Human authority & Healthcare AI literature warns about automation bias, deskilling, liability, and workflow fit \cite{goddard-et-al-2012}; \cite{natali-et-al-2025}; \cite{kelly-2026}; \cite{zheng-et-al-2023}. & The system may teach, score, explain, and log; instructors and supervisors remain responsible for real competency judgments. \\
\hline
Workflow-first learner UI & General usability guidance supports visible status, efficient action paths, error prevention, and recoverable mistakes. Human-AI guidance supports efficient invocation, dismissal, and correction when AI guesses wrong. Clinical alert literature warns that low-value alerts and added tasks can undermine acceptance \cite{amershi-et-al-2019}; \cite{olakotan-and-yusof-2021}; \cite{canovas-segura-et-al-2023}. & Design around the learning mode, use clear mode boundaries, provide fast practice feedback, and suppress hints during assessment. \\
\hline
CV feasibility as an external assumption & Few-shot object detection and surgical-instrument CV papers show that data-efficient detection and constrained instrument recognition are plausible \cite{xin-et-al-2024}; \cite{deol-et-al-2024}; \cite{atabuzzaman-et-al-2025}. & Cite existing papers for feasibility and reserve local CV accuracy for future support-layer validation. \\
\hline
CV generalization as a future validation problem & Kienle et al. show a large cross-manufacturer performance drop despite strong in-domain instrument-stand results \cite{kienle-et-al-2025}. & Any future camera-supported feature needs site-specific validation. The current project can function through manual simulated sorting. \\
\hline
Lookalike tools as learning content & HOSPITools and Atabuzzaman et al. highlight subtle instrument differences; Alfred et al. identifies nomenclature and training as assembly factors \cite{rodrigues-et-al-2022b}; \cite{atabuzzaman-et-al-2025}; \cite{alfred-et-al-2021}. & Teach distinguishing features in retrieval-first cards and track \texttt{misidentified} errors separately from generic wrong answers. \\
\hline
Student participants as novice evidence & Simulation and learning-science evidence supports novice practice studies, while SPD certification and work-system evidence show real work requires supervised hands-on competence \cite{cook-et-al-2011}; \cite{mcgaghie-et-al-2011}; \cite{hspa-crcst-accessed-2026}; \cite{alfred-et-al-2021}. & Student results can support first-use clarity and novice simulated learning. SPD adoption and safety readiness belong in future validation. \\
\hline
Admin-facing metrics and traceability & Work-system studies and traceability reviews support recordkeeping for quality improvement, while economic AI reviews warn that technical performance alone is an incomplete business case \cite{alfred-et-al-2021}; \cite{fayad-et-al-2025}; \cite{vithlani-et-al-2023}; \cite{kastrup-et-al-2024}. & Export learner progress, repeated weak items, time, confidence, and error patterns. Treat ROI as future pilot work. \\
\hline
\end{tabularx}
\end{table}

\section{Implementation Defaults}
\label{sec:appendix-b-implementation-defaults}

\begin{table}[H]
\centering
\caption{Implementation defaults and critique-ready interpretation}
\label{tab:implementation-defaults-and-critique-ready-interpretation}
\small\renewcommand{\arraystretch}{1.2}\setlength{\tabcolsep}{4pt}
\begin{tabularx}{\textwidth}{|>{\raggedright\arraybackslash}X|>{\raggedright\arraybackslash}X|>{\raggedright\arraybackslash}X|}
\hline
Default & Current Use & Critique-Ready Interpretation \\
\hline
File-backed tray and instrument data & Fast local authoring for one seeded module. & Good for a capstone project. A real training program would need role permissions, content review, versioning, and integration decisions. \\
\hline
Pre/post simulated tray sorting & Primary training-effectiveness evidence. & Measures simulated local familiarity and produces paired learning metrics. \\
\hline
Confidence ratings & Learner metacognition and instructor review. & Useful for spotting overconfidence and fragile knowledge. \\
\hline
Existing YOLO/Streamlit tray-check code & Demo code and possible future visual support. & Useful as support-layer infrastructure after separate validation. \\
\hline
Local ignored \texttt{data/}, \texttt{runs/}, and \texttt{weights/} folders & Keeps raw images, training outputs, and weights out of git. & Sensible prototype hygiene. Formal deployment would require access controls, retention policy, audit policy, and privacy review. \\
\hline
\end{tabularx}
\end{table}

\bibliographystyle{plain}
\bibliography{trayguard_refs}

\end{document}
